\documentclass[../main.tex]{subfiles}
\begin{document}
%==========================================TIÊU ĐỀ TẬP========================================
\begin{table}[h]
    \begin{adjustwidth}{-1cm}{}
        \begin{tabular}{>{\centering\arraybackslash}p{6cm}|>{\raggedright\arraybackslash}p{12.5cm}}
            \multirow{5}{6cm}
            {\\[-40pt]
                \flaregothic\fontsize{20pt}{20pt}\selectfont \centering
                \color{eptype}HISTOIRE\color{black}\\
                \freshbot\fontsize{60pt}{60pt}\selectfont
                \color{epnum}32A\\[6pt]
                \cafeteria\fontsize{20pt}{20pt}\selectfont\color{epnum}(D-32A)\color{black}
                \\[18pt]
                \color{black}
            }
             &
            % Dòng thứ nhất: tên tập tiếng Nhật
            {\fotskip\fontsize{18pt}{18pt}\selectfont \ruby{潮}{しお}\ruby{風}{かぜ}の\ruby{記}{き}\ruby{憶}{おく}と\ruby{再}{さい}\ruby{会}{かい}の\ruby{湯}{ゆ}\ruby{気}{け}}
            \\[6pt]
            % Dòng thứ hai: tên tập tiếng Anh
             & {\cafeteria\fontsize{18pt}{18pt}\selectfont Scents of the Sea Breeze, Steam of Reunion}                                                                                                        \\[4pt]
            % Dòng thứ ba: tên tập tiếng Việt
             & {\vnmsans\fontsize{18pt}{18pt}\selectfont Ký ức gió biển và làn hơi tái ngộ}                                                                                                                   \\[8pt]
            % Dòng thứ tư: Nhân vật xuất hiện
             & {\caviar{\fontsize{13pt}{13pt}\selectfont Nhân vật: \emp{kuon2} \emp{chloe} \emp{kaien} \emp{marine} \emp{kaiho2} \emp{suzuno2} \emp{ryuuhou2} \emp{naomi2} \emp{kson} \emp{delu} \emp{game}}}
            \\[8pt]
        \end{tabular}
    \end{adjustwidth}
\end{table}
%=========================================TRUYỆN CHÍNH========================================
\color{black}

\txtbx[2]{
    {\color{vol} \uline{Tóm tắt:}} Vụ việc cô nàng thực tập sinh cá voi sát thủ Sakamata Chloe bỏ bê tắm gội suốt tròn một tháng trời đã kích hoạt một chiến dịch áp giải chưa từng có trong lịch sử. Dưới sự chỉ huy đanh thép của bà Hikawa Kaien và sự hiệp lực của Thuyền trưởng Marine, Chloe bị tống lên xe đưa thẳng về ký túc xá Hikawa để thực hiện cuộc tổng tẩy uế thế kỷ. Giữa những pha giãy giụa dở khóc dở cười trên đường đi và sức mạnh trấn áp tuyệt đối của hai người phụ nữ nhà Hikawa, những kỷ niệm thời thơ ấu đầy nắng gió ven biển mười sáu năm trước giữa Kuon và cô bạn thuở nhỏ bất ngờ ùa về.
}

\stpt{-Histoire 1-}
\stcap{Chuyến xe áp giải và màn thanh tẩy thế kỷ}

\pov{Hikawa Kuon}

Bình thường mỗi khi gia đình có việc cần dùng đến chiếc xe bảy chỗ, mẹ tôi—bà Hikawa Kaien—luôn là người nắm trọn quyền sinh sát sau tay lái. Mẹ lái xe có tiếng là dứt khoát và bạo tay, thừa hưởng trọn vẹn phong thái của một người phụ nữ kinh doanh gáo thép vùng cảng biển.

Thế nhưng chiều tối hôm nay, lần hiếm hoi trong năm, tôi bị đẩy vào ghế lái.

Lý do vô cùng đơn giản: chiếc ghế phụ phía sau đang chứa một ``kiện hàng sinh học nguy hiểm cấp độ ba''. Và ở vị trí áp giải đó, không ai trên cõi đời này đủ uy lực và dũng khí bằng mẹ tôi.

``Kuon, lái êm tay thôi đấy con,'' mẹ tôi vừa gằn giọng vừa siết chặt hai cánh tay bọc găng cao su dầy cộm. ``Đừng có phanh gấp kẻo con cá voi này nó lộn tùng phèo ra sàn xe, dơ hết cả lớp nỉ mẹ mới bọc tuần trước!''

Ở hàng ghế sau, cô nàng thực tập sinh lừng danh của Secret Society holoX, Sakamata Chloe, đang bị trùm kín mít từ đầu đến chân bằng một tấm chăn dã ngoại cách ly chuyên dụng, chỉ thò mỗi cái chỏm đầu tóc xám rối bù và đôi mắt đỏ ngấn lệ ra ngoài. Mẹ tôi ngồi ngay sát bên cạnh, một tay kẹp chặt nách cô nàng, một tay ghì vai như một võ sĩ Judo đang khóa khớp đối thủ trong trận chung kết thế giới.

``P-Poe...! Bác Kaien ơi, cháu ngộp thở chết mất thôi! Bác nới lỏng tay ra một tí cho cháu thở với ạ!'' Chloe mếu máo giãy giụa như một con cá mắc cạn, hai chân quẫy đạp lung tung trong bọc chăn.

``Im ngay!'' Mẹ tôi đanh thép nạt một tiếng khiến kính xe rung rinh. ``Thở ít thôi cho đỡ tỏa cái mùi nồng nặc này ra! Một tháng không tắm mà còn dám đòi hỏi quyền con người hả!? Bác mà thả lỏng tay là cháu lại giở trò nhảy qua cửa sổ chạy trốn chứ gì, bác lạ gì cái nết lươn lẹo của cháu từ hồi bốn tuổi!''

Ngồi ở ghế phụ cạnh tôi, Marine - người đã nhiệt tình hỗ trợ tóm cổ Chloe tại văn phòng và xin đi quá giang một đoạn ra ga tàu điện ngầm - vẫn chưa hết bàng hoàng. Em ấy đeo tới hai lớp khẩu trang y tế, tay cầm chai xịt khử khuẩn liên tục xịt lia lịa vào khoảng không xung quanh lỗ thông gió điều hòa.

Marine chắp tay trước ngực, giọng thảng thốt qua lớp khẩu trang: ``Đội ơn trời Phật, đội ơn bác gái uy quyền tối cao\dots\ Cháu thề là nếu hôm nay bác không kịp thời xuất hiện như một vị cứu tinh, chắc cái văn phòng \hll\ bên đó đã biến thành bãi rác sinh hóa và cháu đã phải nhập viện rửa ruột lần thứ hai rồi!''

Tôi vừa đạp ga cho xe hòa vào dòng xe cộ rời khỏi Tokyo, vừa liếc mắt nhìn qua gương chiếu hậu, khẽ cười trừ: ``Anh không nghĩ là mọi thứ sẽ thảm họa đến cỡ đó đâu. Em cứ nói quá lên\dots''

Chloe nghe thấy thế, lập tức trừng mắt bắn tia nhìn tóe lửa về phía gáy tôi qua gương: ``Tên Kuon chết tiệt kia! Ông còn dám thêm dầu vào lửa à!? Bạn bè thanh mai trúc mã kiểu gì mà bán đứng nhau không chớp mắt thế hả!?''

``Tôi cứu mạng bà đấy chứ,'' tôi điềm nhiên đánh lái qua một khúc cua. ``Bà mà để cái mùi này ám thêm một tuần nữa, khán giả người ta đến bắt đền công ty thì ai chịu trách nhiệm? Chưa kể quản lý của bà vừa nghe mẹ tôi thông báo là đã ký giấy ủy quyền khẩn cấp, còn gửi kèm lời nhắn: `Làm ơn hãy tống nó vào bồn tắm giùm chúng tôi, tiền nước bao nhiêu quản lý chúng tôi xin thanh toán đủ!' đấy thôi.''

Chloe nghẹn họng, hai má đỏ bừng lên vì xấu hổ và tức tối. Chiếc xe lướt êm qua cây cầu vượt nối liền hai khu vực, hướng thẳng về phía ga tàu điện ngầm trung tâm.

Sau khi đỗ xe ở sảnh trả khách của ga, tôi hạ kính cửa phụ xuống. Marine vội vã tháo dây an toàn, trước khi bước xuống tàu còn không quên nghiêng mình cúi đầu thật sâu trước mẹ tôi: ``Cháu cảm ơn bác Kaien nhiều lắm ạ! Chúc bác và Kuon-senpai đại thắng trở về! Làm ơn hãy kỳ cọ cho em ấy thật sạch bong kin kít giùm toàn thể anh chị em công ty nhé bác!''

``Yên tâm đi cháu gái!'' Mẹ tôi vẫy tay đầy tự tin. ``Cứ giao cho bác, đêm nay mà không biến nó thành con cá voi thơm tho như hoa cúc thì bác không mang họ Hikawa!''

``Mẹ cứ thề thốt linh tinh làm gì vật mẹ\dots'' tôi toát mồ hôi hột, cười gượng nhìn nàng Thuyền trưởng dần khuất bóng.

Cánh cửa xe đóng sầm lại. Chiếc xe tiếp tục lăn bánh trên đường cao tốc hướng về phía thành phố Kawachi.

Lúc này, nhận thấy đồng minh duy nhất bên ngoài đã rời đi và chiếc xe đang băng băng về phía ``tử địa'', Chloe bắt đầu rơi vào cơn hoảng loạn tột độ. Cô nàng co rúm hai chân, bắt đầu giở đủ mọi chiêu trò để câu giờ.

``A\dots\ Á!'' Chloe bỗng ôm chặt lấy bụng dưới, mặt mày nhăn nhó biến dạng, cất giọng the the khóc lóc: ``Kuon ơi! Bác Kaien ơi! Cứu cháu với! Cháu\dots\ cháu mắc đi vệ sinh quá rồi!! Bụng cháu nó đang réo ùng ục lên đây này! Dừng xe lại cho cháu xuống trạm xăng đi vệ sinh mau lên, không là nó nổ tung ra xe bây giờ!!''

Mẹ tôi thoáng nhíu mày, hơi ngập ngừng nhìn sang tôi. Dù sao thì chuyện sinh lý khẩn cấp cũng là thứ khó lòng kiểm soát.

Thế nhưng, tôi - người đã cùng lớn lên với cô nàng suốt sáu năm trời ở làng chài ven biển - lại quá đỗi thấu hiểu từng cái nháy mắt, từng cái co chân giả vờ của con bé cá voi sát thủ này. Hồi tám tuổi, mỗi lần trốn mẹ đi tắm biển hay bị bắt quỳ gối học bài, con bé này đều dùng đúng cái bài kêu đau bụng đi vệ sinh này để lẻn cửa sau tẩu thoát.

Tôi không hề giảm tốc độ, mắt vẫn dán thẳng vào con đường phía trước, buông một câu đáp trả dửng dưng như không: ``Thì cứ xả tại chỗ đi. Đằng nào đống quần áo trên người bà tí nữa cũng phải vứt vào lò giặt tẩy trùng, ghế xe thì có bọc chống thấm rồi. Xả xong về nhà tắm một thể, tiện cả đôi đường.''

Mẹ tôi nghe xong, lập tức gật gù tán đồng cái rụp: ``Đúng! Thằng Kuon nói chí phải! Cứ việc tự nhiên đi cháu, bác không chê đâu!''

Chloe nghe xong hai câu nói tàn nhẫn và bình thản đến mức lạnh người ấy thì chết lặng toàn tập.

Đôi mắt đỏ của cô nàng trợn tròn lên hết cỡ, miệng há hốc kinh ngạc như vừa nhìn thấy hai con quái vật vô cảm. Con bé giãy nảy lên, hét toáng qua lớp chăn bông: ``Áaaaa! Hai mẹ con nhà này mất nhân tính vừa thôi chứ!! Tôi là con gái đấy! Là idol triệu sub hẳn hoi đấy nhé!! Ai lại bảo người ta xả thẳng ra quần trên xe bao giờ không hả, đồ độc ác!!''

Từ chiếc đồng hồ thông minh trên cổ tay trái của tôi, giọng nói máy móc của Game cất lên cắt ngang màn than vãn: ``Phân tích sinh trắc học hoàn tất: Nhịp tim của Sakamata-san là một trăm hai mươi nhịp một phút do hoảng sợ, áp suất bàng quang ghi nhận ở mức mười lăm phần trăm, ngưỡng hoàn toàn bình thường. Đây chắc chắn là hành vi ngụy tạo triệu chứng nhằm mục đích đào tẩu. Kuon-user, tiếp tục duy trì tốc độ như vậy đi, không cần phải dừng xe đâu.''

``\textbf{Game, ông im mồm ngay cho tôiiiii!!}'' Chloe gào khóc tuyệt vọng, hoàn toàn buông xuôi gục đầu vào thành ghế.

\ct

Khoảng hơn một tiếng sau, chiếc xe bảy chỗ cuối cùng cũng rẽ vào con dốc quen thuộc của phường Anwa, dừng lại trước cửa ký túc xá Hikawa.

Vừa tắt máy xe, mẹ tôi đã lập tức mở toang cửa sau, lôi xềnh xệch Chloe xuống đất. Thế nhưng vừa nhìn thấy cánh cửa sắt quen thuộc của căn nhà, bản năng sinh tồn trỗi dậy khiến cô nàng cá voi sát thủ vùng vẫy dữ dội. Nhờ cơ thể trơn tuột và sức quẫy của loài thủy tộc, Chloe suýt chút nữa đã tuột khỏi vòng tay của mẹ tôi để lao ra đường cái.

Mẹ tôi vừa túm chặt cổ áo con bé vừa hét lớn về phía cửa nhà: ``\textbf{Kaiho đâu rồi!? Mau ra đây đỡ mẹ một tay! Con cá voi này nó khỏe như trâu mộng ấy, mẹ giữ một mình không xuể!}''

``\textbf{Dạ!}''

Cánh cửa sảnh tầng một bật mở. Kaiho bước ra trong chiếc áo thun ba lỗ và quần thun thể thao gọn gàng. Nghe tiếng mẹ gọi, cô em gái song sinh của tôi chẳng cần hỏi đến câu thứ hai. Bằng phản xạ sắc bén của một thợ cơ khí lành nghề, Kaiho lướt tới như một cơn gió, hai cánh tay rắn chắc như gọng kìm thép vươn ra khóa chặt lấy hai bả vai của Chloe, rồi nhẹ nhàng ghì thẳng cô nàng xuống mặt sàn sảnh gạch đánh *RẦM* một cái nhẹ.

Toàn bộ sức lực giãy giụa của loài cá voi sát thủ lập tức bị triệt tiêu hoàn toàn trước sức mạnh cơ bắp phi thường của nàng Tổng quản \hls.

``Poe\dots!'' Chloe nghẹn thở, hai má dán chặt xuống sàn gạch mát lạnh, mắt nổ đom đóm: ``S-Sức mạnh quái vật gì thế này\dots''

Kaiho khoanh một tay, tay kia vẫn ghì chặt lưng Chloe, nhướng mày nhìn mẹ: ``Kiện hàng bốc mùi mà Aniki báo trước đây hả mẹ?''

``Chính nó đấy!'' mẹ tôi thở phào, quẹt mồ hôi trên trán. ``Xách cổ nó lên phòng tắm tầng ba ngay cho mẹ! Nước nóng pha muối khoáng với xà phòng thảo mộc mẹ đã bật sẵn bình đun từ chiều rồi!''

Đúng lúc đó, từ góc cầu thang, Suzuno khẽ bước xuống. Cô em gái thứ ba của tôi mặc chiếc tạp dề màu be dịu dàng, trên tay cầm sẵn một chiếc thau nhựa lớn và đôi găng tay cao su dài.

Ngửi thấy mùi tanh nồng xộc lên từ người vị khách lạ, em không hề tỏ ra khó chịu hay chê bai. Đôi mắt ánh lên sự chu đáo, ân cần quen thuộc. Suzuno cúi người, nhẹ nhàng gỡ từng lớp áo khoác ngoài và chiếc khăn choàng bẩn thỉu của Chloe cho vào thau nhựa, cất giọng nhỏ nhẹ: ``Mẹ và Onee-sama cứ đưa bạn lên phòng tắm đi ạ. Đống quần áo này cứ để con mang lên sân thượng ngâm dung dịch tẩy rửa sinh học và giặt sấy khử khuẩn cho. Bữa tối con cũng đang hầm sẵn sườn nướng với canh củ quả rồi, tắm xong là vừa kịp ăn ạ.''

Chloe nằm bẹp dưới sàn, ngước nhìn gương mặt thánh thiện như thiên thần của Suzuno, trong một giây ngắn ngủi đã ngỡ mình vừa nhìn thấy đức mẹ cứu rỗi. Con bé yếu ớt thều thào: ``Em gái thiên thần ơi\dots\ cứu chị với\dots''

Nhưng đáp lại lời thỉnh cầu ấy, Suzuno chỉ mỉm cười thật hiền từ, khẽ cúi đầu chào rồi nhanh thoăn thoắt bê thau quần áo bẩn bước lên cầu thang.

``Đi thôi, hết giờ xin xỏ rồi cô nương!''

Kaiho chẳng nói chẳng rằng, một tay nhấc bổng cả người Chloe lên như xách một bao gạo ba mươi cân, sải bước thẳng lên cầu thang dẫn lên phòng tắm tầng ba, theo sau là mẹ tôi với chiếc bàn chải gai và hũ muối tẩy tế bào chết lăm lăm trên tay.

*RẦM!*

Cánh cửa phòng tắm tầng ba đóng sầm lại, chốt khóa đánh *CẠCH* một cái khô khốc.

Và ngay sau đó, những âm thanh kinh hoàng nhất của cuộc ``thanh tẩy thế kỷ'' chính thức vang dội khắp cả tòa nhà Hikawa: tiếng vòi sen xả nước xối xả, tiếng bọt xà phòng tung tóe, tiếng kỳ cọ sột soạt rợn người, hòa cùng tiếng gào khóc thảm thiết, đòi sống đòi chết không ngừng nghỉ của cô nàng cá voi sát thủ.

``\textbf{ÁAAAA! NÓNG QUÁ!! BÁC KAIEN ƠI, DA CHÁU TRÓC HẾT RA RỒI!!}''

``\textbf{Ngồi im! Da cá voi dày thế này tróc thế nào được! Đứng yên cho chị mày kỳ cái lưng xem nào!}''

``\textbf{CỨU TÔI VỚI KUON ƠI!! TÔI THỀ TỪ NAY MỖI NGÀY TẮM BA LẦN MÀ!! ÁAAAAA!!}''

Tôi đứng tựa lưng vào chân cầu thang tầng ba, đưa tay xoa xoa hai bên màng nhĩ đang rung lên bần bật, khẽ thở dài một hơi đầy thương cảm nhưng khóe môi lại không kìm được một nụ cười nhẹ nhõm.

Mười sáu năm rồi, cái cảnh tượng hỗn loạn và ồn ào này mới lại tái diễn. Mùi xà phòng hoa cúc quen thuộc từ tầng ba thoang thoảng bay xuống, mang theo hương vị của những ngày hè ấu thơ đầy nắng gió năm nào.

\stpt{-Histoire 2-}
\stcap{Bữa cơm gia đình và ký ức ven biển}

Trong khi cánh cửa phòng tắm tầng ba vẫn đang rung lên bần bật bởi cuộc vật lộn giữa hai thế lực - một bên là cô nàng cá voi sát thủ giãy chết và một bên là liên quân mẹ Kaien cùng Kaiho, thì tôi đã chọn cho mình một vị trí an toàn hơn nhiều: chiếc ghế sô-pha nỉ dài ở phòng sinh hoạt chung ngay bên đối diện phòng tắm.

Tôi vừa rót cho mình một tách trà lúa mạch mát lạnh từ bình giữ nhiệt thì cánh cửa sảnh tầng ba bỗng bật mở toang.

``\textbf{Chuyện gì xảy ra vậy!? Nhà mình có án mạng à!?}''

Ryuuhou xông vào phòng như một cơn lốc. Con bé vẫn còn ngậm dở một que kem đậu xanh trên miệng, trên tay lăm lăm chiếc gậy bóng chày bằng nhôm sáng loáng, mái tóc tém dựng ngược lên vì hoảng hốt. Ngay sau lưng, bé út Naomi cũng lũn cũn chạy theo, hai tay ôm chặt chú búp bê vải, tròn xoe đôi mắt màu xanh lục bảo nhìn về phía cánh cửa phòng tắm đang bốc khói nghi ngút.

``Bình tĩnh nào Ryuu, hạ cái gậy xuống kẻo đập trúng tivi bây giờ,'' tôi điềm nhiên nhấp một ngụm trà, chỉ tay về phía cánh cửa đang phát ra những tiếng kêu rên: ``Waah! Bác Kaien ơi cháu lạy bác tha cho cái nách của cháu!''. ``Chả có trộm cướp hay án mạng gì đâu. Mẹ với Kaiho chỉ đang làm công tác vệ sinh dịch tễ cho một vị khách đặc biệt thôi.''

Ryuuhou chớp chớp mắt, đưa đầu gậy bóng chày gãi gãi cằm, tai dỏng lên nghe tiếng bọt xà phòng lép bép và tiếng quát đanh thép của cô chị Hai: ``Đứng thẳng cái lưng lên! Lớn tồng ngồng thế này rồi mà để ghét đóng thành từng mảng thế hả!'' Cô nàng tomboy bất giác rùng mình một cái, da gà da vịt nổi rần rần khắp hai cánh tay: ``Khiếp vãi\dots\ Tiếng thét xé lòng nghe rợn cả tóc gáy! Em thề là tiếng lợn bị chọc tiết dưới quê đợt Tết năm ngoái nghe còn êm tai hơn tiếng la hét này! Mà rốt cuộc là ai đang bị tra tấn dã man thế kia hả Onii?''

``Một người quen cũ thôi,'' tôi khẽ mỉm cười.

Đúng lúc đó, tiếng bước chân lộp cộp từ cầu thang tầng trên lại vang lên.

Kson từ tầng bảy thong thả bước xuống, trên người chỉ mặc một chiếc áo ba lỗ rộng thùng thình in chữ rock cỡ lớn và chiếc quần đùi thể thao, tay cầm lon nước ngọt có ga vừa bật nắp đánh *xì*. Ngay bên cạnh cô là Delutaya trong bộ váy mặc nhà màu lam nhạt thanh nhã, trên tay bưng chiếc khay nhỏ đựng mấy đĩa bánh quy bơ vừa nướng xong thơm phức.

``Yo, Kuon-paisen!'' Kson tựa lưng vào khung cửa, ngửa cổ tu một ngụm nước ngọt rồi nhếch mép cười khà khà: ``Trên tầng bảy đóng kín cửa kính mà em vẫn nghe thấy tiếng gầm rú như có con thủy quái Kraken nào bị kéo lên bờ thế này? Nghe đồn bên văn phòng \hll\ vừa tóm được một con hàng khủng bốc mùi đem về đây à?''

Delutaya khẽ đặt đĩa bánh xuống bàn trà, mỉm cười dịu dàng nhìn tôi: ``Em ngửi thấy mùi xà phòng hoa cúc nồng nặc từ tầng một lên tới tận tầng hai luôn đấy, Kuon-kun. Vị khách đó\dots\ là một cô bé trong nhóm tân binh holoX đúng không ạ?''

Suzuno vừa lúc từ sân thượng phơi đồ đi xuống, tay cầm chiếc khăn lau tay sạch sẽ, cũng nhẹ nhàng kéo ghế ngồi xuống cạnh Naomi. Cả gia đình—từ các em gái cho đến hai cô nàng cư dân đặc biệt của các tầng trên—đều quây quần đông đủ quanh chiếc bàn gỗ, ánh mắt ai nấy đều dồn về phía tôi với sự tò mò tột độ.

Tôi đặt chén trà xuống bàn, ngả người ra sau lưng ghế sô-pha, ánh mắt bất giác nhìn về phía làn hơi nước ấm áp đang lọt qua khe cửa phòng tắm. Những ký ức vốn đã ngủ yên suốt mười sáu năm qua bỗng chốc như một thước phim quay chậm, sống động ùa về trong tâm trí.

``Chuyện này\dots\ thực ra phải bắt đầu từ cuộc trò chuyện tối hôm qua ở bàn ăn nhà mình cơ,'' tôi chậm rãi cất lời.

\begin{center}
    \uline{Hồi tưởng.}
\end{center}

Buổi tối, một tuần trước ngày Chloe chính thức giáp mặt tôi ở công ty.

Bữa cơm tối hôm ấy tại căn bếp tầng năm diễn ra trong một không gian yên ắng lạ thường—bởi trên bàn ăn lúc này chỉ có vỏn vẹn ba người: bố mẹ và tôi.

Các em gái đều vắng nhà vì những lý do riêng: Kaiho và Suzuno bận dắt bé Naomi sang khu trung tâm thương mại quận Kyuukami để mua sắm thêm đồ dùng học tập và vài bộ quần áo mùa hè mới, tiện thể cho Naomi đi ăn kem và ghé nhà sách; trong khi cô nhóc Ryuuhou thì đã xin phép từ chiều để ở lại trường tham gia buổi tập dượt tổng duyệt ban nhạc chuẩn bị cho đêm nhạc giao lưu cuối tuần.

Căn bếp chỉ còn lại ba chiếc bát, ánh đèn vàng ấm áp rọi xuống mâm cơm gia đình giản dị nhưng đậm đà. Bố ngồi tựa lưng vào ghế gỗ, chậm rãi nhấp một ngụm rượu sake cay nồng sau một ngày dài miệt mài bên máy phay máy tiện dưới xưởng. Mẹ vừa xới thêm cho tôi một bát cơm đầy ăm ắp vừa mỉm cười thở phào: ``Hiếm hoi lắm mấy đứa nhỏ mới đi vắng hết một lượt, nhà cửa tự dưng rộng thênh thang, ăn bữa cơm mà thấy tĩnh tâm hẳn ra.''

Tôi mỉm cười gắp một miếng cá kho đậm vị: ``Nhà đông con thì ồn ào quen rồi, tự nhiên vắng tụi nó con lại thấy thiêu thiếu tiếng chành chọe của Ryuu với Kaiho.''

Đúng lúc ấy, bố vừa lướt ngón tay trên màn hình chiếc máy tính bảng cũ để đọc mục tin tức giải trí buổi tối, bỗng khẽ nheo mắt lại, nhấp thêm một ngụm rượu rồi ngẩng lên nhìn tôi: ``Kuon này. Dạo gần đây công ty con vừa cho ra mắt một nhóm thực tập sinh mới tên là holoX gì đó rầm rộ lắm đúng không? Bố thấy báo mạng người ta đưa tin liên tục, bảo là nhóm tân binh thế hệ thứ sáu có phong cách rất độc lạ.''

``Vâng ạ,'' tôi đáp, vừa nuốt miếng cơm. ``Nhóm gồm năm thành viên, lấy hình tượng một tổ chức bí mật mưu đồ xâm chiếm thế giới. Bên truyền thông và cộng đồng mạng đang chú ý dữ dội lắm bố ạ. Lịch trình tập luyện và ghi hình của họ dày đặc suốt cả tháng nay.''

``Trong nhóm năm đứa đó\dots\ hình như có một đứa bé mang họ Sakamata, tên là Chloe đúng không con?'' bố nheo mắt, dí sát màn hình máy tính bảng lại gần chiếc kính lão để soi một bức ảnh chụp tập thể trong buổi họp báo ra mắt.

Cái tên vừa thốt ra khỏi miệng bố, đôi đũa trên tay mẹ tôi bỗng nhiên khựng lại giữa không trung.

Mẹ quay ngoắt người lại, hai mắt mở to đầy kinh ngạc: ``Sakamata Chloe!? Anh vừa bảo con bé tên là Sakamata Chloe thật đấy à!?''

Mẹ vội vàng buông đũa, giật phắt chiếc máy tính bảng từ tay bố, kéo kính lão xuống tận sống mũi để soi thật kỹ từng đường nét của cô gái trùm chiếc áo choàng xám mang mặt nạ cá voi sát thủ trên màn hình. Chỉ sau đúng vài giây quan sát đôi mắt đỏ đặc trưng, chiếc răng khểnh và nụ cười lém lỉnh quen thuộc ấy, mẹ đã vỗ đùi đánh *ĐÉT* một cái vang dội cả gian phòng: ``Trời ơi! Đúng là nó rồi! Đúng là con bé Chloe nhà Sakamata sát vách hàng rào nhà mình hồi ở làng chài ven biển đây mà! Thảo nào hồi đầu năm, bà Sakamata có gọi điện tâm sự với mẹ cả tiếng đồng hồ, bảo con bé Chloe một mực đòi rời thị trấn ven biển khăn gói lên Tokyo để theo đuổi con đường giải trí streamer gì đó! Mẹ con bà ấy còn cãi nhau một trận to vì mẹ nó lo con gái thân cô thế cô nơi đất khách quê người. Thế mà giờ nó trúng tuyển thẳng vào công ty của thằng Kuon thật này!''

Bố đặt ly rượu xuống bàn, gương mặt sương gió dạn dày của người thợ cơ khí khẽ giãn ra một nụ cười ấm áp, hoài niệm: ``Con bé cá voi nhỏ xíu đen nhẻm hay bám đuôi thằng Kuon ngày xưa đấy à? Thoắt cái mười sáu năm trôi qua, giờ đã lớn tướng thành idol nổi tiếng triệu người theo dõi rồi cơ đấy. Đúng là thời gian trôi nhanh như bóng câu qua cửa sổ.''

Mẹ buông chiếc máy tính bảng xuống bàn, ánh mắt nhìn tôi bỗng chốc trở nên nghiêm cẩn và chất chứa đầy sự lo lắng, trăn trở của một người bề trên: ``Kuon này, con nghe mẹ dặn đây. Hai đứa con lớn lên cùng nhau từ thuở cởi truồng tắm mưa, con là người hiểu tính nết nó nhất. Con bé Chloe bề ngoài thì mồm mép tép nhảy, lúc nào cũng tỏ ra bất cần, lanh chanh tinh tướng thế thôi, chứ thực chất bên trong nó là đứa cực kỳ hướng nội, thiếu cảm giác an toàn và rất dễ tủi thân. Rời quê lên chốn Tokyo phồn hoa đô hội này một thân một mình, không người thân thích, chắc chắn con bé phải chịu đựng biết bao nhiêu áp lực, cô đơn và sợ hãi mà không dám hé răng tâm sự với ai đâu.''

Bố gật đầu tiếp lời, giọng trầm ấm và đầy sức nặng: ``Con là anh lớn, lại giữ vị trí Tổng quản lý bên đó. Ở công ty, nếu con bé có bướng bỉnh, quậy phá hay làm loạn, con cứ nghiêm khắc uốn nắn theo kỷ luật công việc. Nhưng nhớ lời bố mẹ: Về mặt đời sống, con phải luôn để mắt bảo bọc, chở che và chăm sóc bạn giùm bố mẹ. Đừng để con bé cảm thấy bơ vơ, lạc lõng giữa đất khách quê người nghe chưa con. Bạn bè thuở ấu thơ tìm lại được nhau giữa biển người mênh mông này là cái duyên lớn lắm đấy.''

``Con nhớ rồi ạ,'' tôi gật đầu, trong lòng dâng lên một cảm giác ấm áp khó tả.

Thế nhưng ngay sau đó, mẹ Kaien bỗng thở dài thườn thượt, lắc đầu chép miệng bồi thêm một câu đầy vẻ bất an: ``Nhưng mẹ lo nhất là cái nết lười tắm của nó! Hồi nhỏ nó sợ nước nóng với ghét tắm có tiếng khắp cả cái làng chài, mỗi lần bị bắt đi tắm là giãy đành đạch như cá mắc cạn. Giờ lên thành phố sống một mình, bận rộn tập luyện thâu đêm suốt sáng, khéo nó lại tha nguyên cái tật ở dơ ấy đi làm thì khổ cả công ty con ạ! Nếu thấy nó bốc mùi thì con phải lôi cổ nó về đây ngay nghe chưa!''

Lúc ấy, tôi chỉ cười trừ nghĩ mẹ lo xa. Ai mà ngờ được rằng lời tiên tri đầy ngao ngán của mẹ tối hôm đó lại linh nghiệm đến mức rợn người chỉ sau chưa đầy một tuần sau.

\begin{center}
    \uline{Kết thúc hồi tưởng.}
\end{center}

Nghe tôi thuật lại cặn kẽ cuộc trò chuyện tối qua của bố mẹ, Kson ngồi cạnh bật cười sặc sụa, vỗ đùi bôm bốp: ``Haha! Hóa ra bác gái đã dự đoán trước được tương lai từ tối hôm qua rồi à!? Đúng là gừng càng già càng cay! Bác ấy đã lên sẵn chiến dịch thanh tẩy từ trước khi con cá voi này kịp trở tay!''

Delutaya mỉm cười tủm tỉm, khẽ cắn một miếng bánh quy bơ: ``Nghe bố mẹ anh Kuon nói chuyện ấm áp thật đấy\dots\ Hóa ra cô bé ấy là bạn thanh mai trúc mã của Kuon-kun từ thời thơ ấu ven biển sao?''

Ryuuhou thì chống cằm lên mặt bàn, mắt lấp la lấp lánh vẻ tò mò không kìm được: ``Này Onii, thế hồi bốn tuổi ở quê hai người gặp nhau thế nào? Em nghe bảo hồi xưa anh lười vận động lắm, chỉ thích ngồi ườn một chỗ, thế mà lại đi chơi thân sống chết được với một con cá voi tăng động á? Kể chi tiết từ đầu xem nào!''

Tôi dựa lưng vào ghế sô-pha, ánh mắt nhìn xa xăm về phía trần nhà, giọng nói trầm xuống đưa mọi người quay ngược bánh xe thời gian về mười sáu năm trước.

\begin{center}
    \uline{Hồi tưởng mười sáu năm trước.}
\end{center}

Đó là mùa hè năm tôi vừa tròn bốn tuổi, khi gia đình tôi còn sinh sống tại một thị trấn làng chài ven biển thanh bình quanh năm lộng gió mặn.

Nhà Hikawa và nhà Sakamata nằm sát vách nhau, chỉ ngăn cách bởi một hàng rào gỗ bạch đàn thấp lè tè, thưa thớt đến mức một đứa trẻ bốn tuổi chỉ cần nghiêng mình là có thể chui tọt qua dễ dàng. Bố mẹ hai nhà vốn là chỗ quen biết lâu năm trong xóm chài, nhưng mãi đến buổi chiều oi ả ấy, tôi và con bé mới có cuộc chạm trán định mệnh đầu tiên.

Hôm ấy, trời nắng như đổ lửa. Tôi - một thằng nhóc bốn tuổi lười biếng, sợ ánh nắng chói chang - đang ngồi co ro dưới bóng râm mát rượi của tán cây bàng trước sân nhà, say sưa lật giở từng trang của cuốn truyện tranh bách khoa động vật cũ mèm.

Đột nhiên, từ bờ rào gỗ bên cạnh phát ra một tiếng *RẮC!* gãy gọn.

Tôi giật mình ngẩng đầu lên.

Bám trên thanh gỗ rào là một cô bé nhỏ thó, gầy nhom, mặc chiếc áo phông cộc tay màu xám rộng thùng thình lấm lem vệt kem dâu. Hai bím tóc ngắn ngủn màu xám khói của con bé vểnh ngược sang hai bên, miệng đang ngậm chặt một cây kẹo mút đỏ chót. Trên mái tóc rối bù ấy là một chiếc mũ vải hình đầu cá voi sát thủ kỳ quặc đang đội lệch sang một bên.

Con bé chớp chớp đôi mắt màu đỏ thẫm to tròn nhìn tôi chằm chằm, rồi chẳng chút e dè hay ngượng ngùng, nó nhảy tót qua hàng rào đánh uỵch một cái xuống bãi cát sân nhà tôi, tiến thẳng lại gần: ``Này! Cậu là Hikawa Kuon đúng không? Con trai bác Ryujin thợ sắt đúng không?''

Tôi trố mắt nhìn sinh vật lạ lẫm vừa từ trên trời rơi xuống, ngập ngừng gật đầu: ``Ừ\dots\ Còn cậu là ai?''

Con bé rút que kẹo mút ra khỏi miệng, quẹt ngang mu bàn tay áo đầy vẻ tự hào: ``Tớ là Sakamata Chloe! Cá voi sát thủ vĩ đại nhất của cái bãi biển đằng kia! Mẹ tớ bảo nhà bên này có một đứa bằng tuổi tớ, nhưng sao trưa nắng đẹp thế này mà cậu cứ ngồi thu lu một xó như ông cụ non thế hả? Ra biển chơi với tớ nào!''

``Không đi đâu,'' tôi lắc đầu nguầy nguậy, ôm chặt cuốn truyện tranh vào ngực. ``Ngoài biển nắng lắm, cát nóng rát cả chân, lại còn đầy gió mặn bẩn thỉu nữa. Ngồi đây đọc sách mát hơn nhiều.''

Chloe trừng mắt nhìn tôi như thể vừa nghe thấy một điều sỉ nhục to lớn đối với đại dương. Con bé chẳng thèm nói thêm lời thứ hai, lập tức lao tới túm chặt lấy cổ tay tôi bằng một sức kéo dẻo dai đến kinh ngạc của loài thủy tộc: ``Đọc sách cái gì mà đọc sách! Trẻ con là phải ra biển! Đi mau lên, tớ vừa phát hiện ra một bãi đá ngầm có bao nhiêu là còng gió với vỏ sò lấp lánh kìa!''

Thế là, bất chấp những tiếng la oai oái và sự kháng cự yếu ớt của một cậu nhóc lười vận động, tôi bị con bé lôi xềnh xệch qua rặng phi lao, hướng thẳng ra bờ biển rực rỡ nắng vàng.

Và cái buổi chiều định mệnh ấy đã mở màn cho một chuỗi ngày ``ác mộng đầy ắp tiếng cười'' kéo dài suốt sáu năm tuổi thơ của tôi.

Kể từ ngày hôm đó, tôi hoàn toàn mất đi sự yên bình. Cứ hễ đến đầu giờ chiều, khi người lớn trong xóm chài bắt đầu chợp mắt, tiếng huýt sáo quen thuộc cùng cái đầu bù xù của Chloe lại nhô lên sau hàng rào gỗ.

``Kuon ơiii! Cầm xô với xẻng ra mau lên! Hôm nay thủy triều rút cạn rồi kìa!!''

Con bé Chloe nhỏ con là thế nhưng lại mang một nguồn năng lượng vô tận và một sự nghịch ngợm không ai bì kịp.

Có những trưa hè ngập nắng, hai đứa dắt díu nhau ra bãi bùn lầy ven cửa sông để cào ngao và bắt còng gió. Chloe xắn ống quần lên tận bẹn, lao thẳng xuống lớp bùn đen quánh như một con rái cá thực thụ. Con bé thoăn thoắt thò tay vào từng hang cua, bất chấp bùn bắn tung tóe từ đầu đến chân, thi thoảng bị một con còng gió kẹp chặt vào đầu ngón tay đau điếng liền nhảy cẫng lên la oai oái: ``Ái ui đau quá! Kuon ơi cứu tớ! Con quái vật này nó dám cắn cá voi sát thủ này!!''

Tôi vừa bò trên cát kéo chiếc xô sắt vừa thở dài: ``Ai bảo cậu thò tay bừa bãi vào hang đá làm gì, dùng cái que gỗ mà khều nó ra chứ! Đứng yên đấy mình gỡ cho!''

Gỡ được con còng ra, ngón tay chảy máu rỉ rỉ, thế mà chỉ cần ba giây sau con bé đã lại toe toét cười, giơ chiến lợi phẩm lên khoe ầm ĩ khắp bãi biển.

Lại có những buổi chiều sóng êm, hai đứa thi nhau ném đá nhảy trên mặt nước. Chloe nhặt những viên đá dẹt mài nhẵn thó, khom người vung tay thật mạnh. Viên đá nảy trên mặt sóng tới năm, sáu nhịp làm con bé nhảy cẫng lên hò reo sung sướng: ``Thấy chưa Kuon! Bảy nhịp đấy nhé! Tớ là bá chủ mặt biển rồi! Cậu mà ném thua là phải cõng tớ về tận cổng đấy nghe chưa!''

``Cậu ăn gian thì có, nhịp thứ sáu là nó chìm nghỉm rồi còn đâu! Với lại người cậu toàn bùn tanh rình thế kia, ai thèm cõng cho bẩn áo!''

``Không biết đâu! Rõ ràng là bảy nhịp! Mau cõng tớ về không tớ mách bác Kaien là cậu xô tớ ngã xuống bùn bây giờ!!''

Những cuộc cãi vã trẻ con vô tư lự ấy cứ thế trôi đi trong tiếng sóng vỗ rì rào và mùi gió biển mằn mặn ngấm vào da thịt. Đến khi hoàng hôn đỏ rực buông xuống đường chân trời, hai đứa mới lếch thếch dắt tay nhau về xóm, mặt mũi đứa nào đứa nấy đen nhẻm, tóc bết chặt cát biển, quần áo ướt sũng và nồng nặc mùi tanh tưởi của bùn lầy.

Thế nhưng, sau niềm vui rực rỡ ngoài biển khơi luôn là một ``cực hình'' khủng khiếp đang chực chờ nơi ngưỡng cửa.

Vừa bước chân đến đầu ngõ, hai đứa đã thấy mẹ tôi và mẹ của Chloe đứng sẵn ở cổng với hai chiếc roi dâu và hai chiếc khăn bông to tướng.

Mẹ tôi xách tai tôi lôi thẳng ra bể nước mưa góc sân: ``Mày nhìn lại cái người mày xem có khác gì con ma bùn không hả Kuon!? Lại trốn đi cào ngao với con Chloe đúng không!? Cởi áo ra mau!''

Trong khi đó, bên kia hàng rào, mẹ của Chloe đã túm chặt cổ áo con gái lôi xềnh xệch vào phòng tắm. Và ngay lập tức, cái bản tính sợ nước nóng và ghét tắm kinh hoàng của loài cá voi sát thủ bắt đầu phát huy tác dụng.

Cứ mỗi lần mẹ nó múc một gáo nước ấm dội từ trên đầu xuống là con bé lại giãy nảy lên như đỉa phải vôi, hai tay bám chặt lấy bậu cửa gào khóc thảm thiết như bị xẻ thịt: ``Oa oa oa! Mẹ ơi tha cho con! Nước nóng nó làm chảy mỡ cá voi của con mất thôi!! Con không tắm đâu! Con thề ngày mai con tự sạch mà!! Áaaaa!!''

Bên này bờ rào, tôi vừa đứng dội nước lạnh vừa ngoẹo cổ nhìn sang, cười đắc chí trêu với sang: ``Đáng đời con cá voi ở bẩn nhé! Cho kỳ sạch lớp ghét trên người đi không người ta tưởng tảng đá rêu phong bây giờ!''

``\textbf{Kuon chết tiệt!! Tớ thù cậuuuuu!!}'' tiếng thét căm hờn của Chloe hòa cùng tiếng bọt xà phòng lép bép đã trở thành một phần âm thanh quen thuộc của xóm chài nhỏ suốt bao năm tháng êm đềm.

Thế nhưng, những ngày hè rực rỡ ấy không thể kéo dài mãi mãi.

Năm tôi mười tuổi, công việc cơ khí của bố tôi có một bước ngoặt lớn. Bố quyết định gom toàn bộ vốn liếng đưa cả gia đình rời khỏi làng chài ven biển, chuyển lên thủ đô Kawachi để mở xưởng cơ khí riêng và lập nghiệp lâu dài.

Tin tức gia đình Hikawa sắp chuyển đi nổ ra như một tiếng sét đánh ngang tai cô bạn nhỏ.

Suốt ba ngày trước khi chuyển nhà, Chloe không hề sang gọi tôi ra biển một lần nào. Con bé nhốt mình trong phòng, thậm chí bỏ cả ăn cơm.

Cho đến buổi chiều muộn của ngày chia tay.

Trời đổ một cơn mưa rào mùa hạ xối xả, gió biển thổi rát rạt qua rặng phi lao. Chiếc xe khách chở gia đình tôi và đống đồ đạc cơ khí nặng trịch đang nổ máy chờ ở bến xe buýt ven thị trấn. Bố mẹ tôi đang bận rộn chất những chiếc hòm gỗ cuối cùng lên nóc xe.

Tôi đứng dưới mái hiên bến xe ngắm nhìn màn mưa trắng xóa, lòng trĩu nặng một nỗi buồn man mác. Đúng lúc ấy, một bóng hình nhỏ nhắn ướt sũng bất ngờ lao xuyên qua làn mưa bão, chạy thục mạng về phía tôi.

Là Chloe.

Con bé không mang ô, chiếc áo phông sũng nước dán chặt vào người, mái tóc xám bết bát nước mưa chảy ròng ròng xuống đôi má nhợt nhạt. Đôi mắt đỏ thẫm của con bé sưng húp, đỏ hoe vì khóc suốt mấy ngày liền.

Con bé chạy đến trước mặt tôi, thở hổn hển rồi bất thình lình đưa hai bàn tay nhỏ bé, run rẩy nắm chặt lấy vạt áo sơ-mi của tôi, siết chặt đến mức các đầu ngón tay trắng bệch ra: ``Kuon\dots\ Cậu\dots\ cậu đi thật đấy à?''

Tôi cúi đầu, cắn nhẹ môi dưới, khẽ gật đầu: ``Ừ. Bố mình bảo lên Kawachi mở xưởng mới có tương lai. Chiều nay xe chạy rồi.''

Nước mắt của Chloe lại trào ra, hòa lẫn vào những giọt nước mưa chảy tràn trên khóe môi. Con bé mếu máo, bờ vai gầy run lên bần bật từng cơn nấc nghẹn ngào: ``Cậu nói dối\dots\ Cậu bảo lớn lên hai đứa mình sẽ cùng nhau đóng một con tàu gỗ ra khơi bắt cá voi khổng lồ cơ mà\dots\ Sao tự nhiên cậu lại bỏ đi một mình lên thành phố chứ\dots''

``Mình không nói dối đâu,'' tôi luống cuống rút chiếc khăn tay trong túi áo ra lau vệt nước mưa trên trán con bé. ``Nhưng bố mẹ chuyển đi thì mình phải đi theo chứ. Sau này lớn lên mình sẽ về thăm cậu mà.''

Chloe lắc đầu quầy quậy, hai tay vẫn siết chặt lấy vạt áo tôi như thể chỉ cần buông tay ra là tôi sẽ tan biến vào hư không mãi mãi: ``Cậu lên thành phố phồn hoa nhiều đồ chơi đẹp, rồi cậu sẽ quên tớ cho mà xem! Cậu sẽ có bạn mới, rồi cậu sẽ chê con bé cá voi làng chài lấm lem bùn đất này cho mà xem!!''

Tôi nhìn cô bạn nhỏ yếu đuối và đầy bất an trước mặt, bất giác đưa tay nắm lấy bàn tay lạnh ngắt của con bé, cất giọng thật rành rọt và quả quyết: ``Không bao giờ có chuyện đó đâu, ngốc ạ! Mình thề đấy. Dù có lên Kawachi, dù có đi đâu xa chăng nữa, mình cũng tuyệt đối không bao giờ quên Sakamata Chloe, con cá voi ngốc nghếch lười tắm nhất cái làng chài này đâu!''

Chloe sụt sịt mũi, ngước đôi mắt đỏ hoe lên nhìn tôi. Đoạn, con bé thò bàn tay run rẩy vào túi quần ướt sũng, lôi ra một vật nhỏ xíu đặt dúi vào lòng bàn tay tôi.

Đó là một chiếc móc khóa nhỏ bằng gỗ trôi dạt dạt vào bờ biển, được chính tay con bé dùng dao gọt vụng về thành hình một chiếc đuôi cá voi nhỏ, bên trên còn khắc nguệch ngoạc chữ cái `C' và `K' lồng vào nhau.

``Cầm lấy cái này đi\dots'' Chloe lí nhí qua tiếng nấc. ``Đây là bùa hộ mệnh của tớ đấy. Cậu phải giữ nó suốt đời nghe chưa. Nếu dám làm mất\dots\ sau này gặp lại tớ sẽ đấm cho cậu bẹp dí đấy!''

Tiếng còi xe buýt vang lên giục giã. Mẹ tôi từ cửa xe ngoái ra vẫy tay: ``Kuon ơi! Vào xe mau lên con, xe chạy bây giờ!''

Tôi siết chặt chiếc móc khóa đuôi cá voi trong lòng bàn tay, khẽ vỗ nhẹ vai Chloe một cái rồi quay người bước lên bậc cửa xe buýt.

Cánh cửa xe khép lại. Chiếc xe từ từ lăn bánh trên con đường nhựa ven biển mờ mịt hơi sương.

Qua ô cửa kính nhòe nhoẹt nước mưa, tôi nhìn lại phía sau. Bóng hình nhỏ nhắn của Chloe vẫn đứng chôn chân dưới mái hiên bến xe buýt cũ kỹ, hai tay giơ cao vẫy lia lịa trong màn mưa gió, miệng không ngừng hét lớn điều gì đó mà tiếng động cơ xe đã át đi mất.

``\textbf{Kuon ơiii! Nhớ giữ lời hứa đấy nhéeee!! Nhất định phải gặp lại nhau đấy!!}''

Hình ảnh cô bé cá voi nhỏ bé đứng cô độc giữa cơn mưa rào làng chài năm ấy đã khắc sâu vào tâm trí tôi suốt cả một thời niên thiếu.

\begin{center}
    \uline{Kết thúc hồi tưởng.}
\end{center}

Tôi dừng lại, khẽ thở dài một hơi nhẹ tênh.

Cả căn phòng sinh hoạt chung ở tầng ba lúc này chìm vào một khoảng lặng sâu lắng đến nghẹt thở. Không còn tiếng cười đùa hay châm chọc thường ngày nữa.

Ryuuhou ngồi bất động bên cạnh tôi, que kem đậu xanh trên tay đã tan chảy từ lúc nào mà không hề hay biết. Đôi mắt sắc sảo của cô nàng tomboy mở to, tràn ngập sự xúc động và kinh ngạc: ``Trời đất ơi\dots\ Câu chuyện chia tay bến xe buýt trong mưa nghe như phim điện ảnh thế này á!? Hóa ra lúc hai người chia tay đẫm nước mắt ở làng chài thì em vẫn còn đang nằm ngoan ngoãn trong bụng mẹ thật! Sau khi chuyển lên Kawachi được đúng một tháng thì em mới oe oe chào đời\dots\ Thảo nào bà chị cá voi ấy nhìn thấy em lại hoang mang ngơ ngác đến thế!''

Bé út Naomi ngồi dưới sàn, hai bàn tay nhỏ xíu siết chặt lấy gấu váy, đôi mắt xanh lục bảo rưng rưng ngấn lệ vì thương cảm: ``Chị Chloe lúc đó chắc là khóc nhiều lắm\dots\ Một mình ở lại bãi biển không có Onii-chan chơi cùng chắc buồn lắm ạ\dots''

Suzuno dịu dàng ôm lấy Naomi vào lòng, đôi mắt hai màu ánh lên sự đồng cảm sâu sắc của một người từng trải qua bao biến cố chia ly: ``Một đứa trẻ hướng nội và sợ cô đơn như em ấy, khi mất đi người bạn thân thiết duy nhất thuở ấu thơ, chắc chắn vết thương lòng ấy lớn lắm. Mười sáu năm qua, chắc em ấy đã giữ mãi hình bóng của anh Kuon trong tim mình.''

Kson ngả người ra sau ghế, im lặng mất vài giây rồi khẽ nhấp một ngụm nước ngọt, cất giọng trầm tư hiếm thấy: ``Mười sáu năm bặt vô âm tín\dots\ Từ một cô bé làng chài lấm lem bùn đất, nuốt nước mắt một mình bước lên Tokyo hoa lệ, chịu đựng biết bao cay đắng của nghề streamer để trở thành một thần tượng nổi tiếng\dots\ rồi cuối cùng lại tình cờ bước vào đúng công ty mà cậu làm Tổng quản lý. Kuon à, trên đời này không có sự trùng hợp nào ngẫu nhiên đến thế đâu. Đó là duyên phận đấy.''

Delutaya khẽ gật đầu, ánh mắt đong đầy sự thấu hiểu của một người từng đi qua giông bão: ``Em nghĩ cô bé ấy hẳn phải xúc động và hạnh phúc lắm. Miệng thì gào thét cãi vã thế thôi, nhưng được gặp lại người bạn thuở ấu thơ, được mẹ của bạn tắm rửa kỳ cọ và được đón nhận vào một gia đình ấm áp như thế này\dots\ chắc chắn mọi sự cô đơn, tủi hờn suốt mười sáu năm qua trong lòng em ấy đang tan biến rồi.''

Đúng lúc câu nói của Delutaya vừa dứt\dots

*CẠCH!*

Tiếng chốt khóa phòng tắm tầng ba bỗng bật mở.

Một làn hơi nước nóng hổi mang theo hương thơm ngạt ngào của xà phòng hoa cúc và tinh dầu tràm cuồn cuộn tỏa ra hành lang như một làn sương mờ ảo. Cả căn phòng sinh hoạt chung bỗng chốc im bặt, mọi ánh mắt đồng loạt quay phắt lại nhìn về phía cánh cửa vừa mở.

Cuộc ``thanh tẩy thế kỷ'' cuối cùng cũng đã đi đến hồi kết.

\stpt{-Histoire 4-}
\stcap{Chloe ``thơm tho'' và những mối duyên tao ngộ}

Từ giữa làn sương mờ ảo cuồn cuộn tràn ra từ cửa phòng tắm, bóng dáng của hai ``vị tướng lĩnh chiến thắng'' lừng lững bước ra trước tiên.

Mẹ tôi xắn cao tay áo quá khuỷu, mồ hôi lấm tấm đọng trên vầng trán cương nghị nhưng đôi mắt lại sáng rực niềm tự hào khôn tả. Đôi găng tay cao su màu cam được mẹ tháo ra vắt hờ bên hông tạp dề, tay phải vẫn còn lăm lăm chiếc bàn chải kỳ lưng cán dài chẳng khác nào một vị hiệp sĩ vừa rút gươm khải hoàn sau trận tử chiến với rồng thiêng.

Ngay bên cạnh, Kaiho vừa bước ra vừa khẽ lắc lắc cổ tay cho đỡ mỏi, chiếc áo thun thể thao lấm tấm vài vệt bọt xà phòng trắng muốt, khẽ thở hắt ra một hơi: ``Xong rồi đấy. Đè con cá voi này nằm yên đúng là tốn sức gấp mười lần nâng tạ một trăm cân ở phòng gym.''

Và rồi, sau lưng hai người họ, ``nạn nhân'' của cuộc tổng vệ sinh dịch tễ cuối cùng cũng ngập ngừng ló dạng.

Mọi người có mặt trong phòng sinh hoạt chung đồng loạt nín thở, tròn xoe mắt nhìn.

Đứng trước mặt chúng tôi lúc này không còn là một cô nàng sát thủ dọn dẹp hôi hám, nhếch nhác với chiếc áo khoác bốc mùi chua nồng và mái tóc bết dính thường ngày nữa.

Chloe bước ra với những bước chân lí nhí, rụt rè chưa từng thấy. Quanh người cô nàng quấn chặt một chiếc khăn tắm bông dày dặn màu trắng muốt to sụ, mép khăn ép chặt vào khuôn ngực nhỏ nhắn và che kín quá nửa đùi, để lộ ra đôi bờ vai thon thả trắng ngần hơi ửng hồng lên vì ngâm mình quá lâu trong làn nước nóng. Mái tóc xám tro mềm mại đã được vắt khô, quấn gọn ghẽ bên trong một chiếc khăn bông thứ hai trên đầu như một chiếc mũ vải tròn xoe, chỉ chừa lại vài lọn tóc tơ còn ẩm ướt dính hờ hững bên gò má.

Hai bàn tay nhỏ nhắn của cô nàng túm chặt lấy vạt khăn trước ngực, đôi mắt đỏ ruby long lanh ngấn nước ngập ngừng nhìn xuống sàn nhà bằng gỗ, hai má đỏ bừng lên như hai trái cà chua chín mọng vì xấu hổ tột cùng.

Và điều kỳ diệu nhất: toàn bộ cái thứ mùi hương ám ảnh được ví như ``vũ khí sinh học cấp độ ba'' kia đã hoàn toàn bốc hơi không còn một dấu vết. Thay vào đó, một làn hương hoa cúc dại thanh khiết, ngào ngạt hòa quyện cùng mùi tinh dầu tràm ấm áp, sảng khoái lan tỏa ngập tràn khắp căn phòng sinh hoạt chung, khiến bất kỳ ai hít vào một hơi cũng cảm thấy sảng khoái đến từng tế bào phổi. Da dẻ cô nàng lúc này sạch bong kin kít, mịn màng và trắng trẻo đến mức phát sáng dưới ánh đèn neon ấm áp của tầng ba.

``P-Poe\dots\ Nhìn cái gì mà nhìn ghê vậy\dots?'' Chloe lí nhí trong cổ họng, hai hàm răng khẽ va vào nhau lập cập, ngượng ngùng bước tới rồi ngồi thu lu một góc trên chiếc ghế sô-pha nỉ, kéo hai đầu gối sát vào ngực như muốn thu nhỏ cơ thể lại hết mức có thể.

Ryuuhou vẫn đang ngậm dở que kem đậu xanh trên miệng và lăm lăm chiếc gậy bóng chày trên tay liền tròn xoe mắt, bước tới gần từng bước một. Con bé khịt khịt mũi hai cái, vẻ mặt từ kinh ngạc tột độ chuyển sang bàng hoàng không tin nổi vào mắt mình: ``Uầy\dots\ Thơm phức luôn kìa! Mẹ với Kaiho-nee dùng phép thuật biến hình đấy à?''

Con bé cúi sát mặt lại ngắm nghía. Đúng lúc ấy, Chloe hơi ngẩng đầu lên. Hai ánh mắt vừa chạm nhau, đôi đồng tử đỏ rực của cô nàng cá voi sát thủ bỗng chốc trợn tròn lên hết cỡ như bị điện giật.

Cô nàng há hốc miệng, chỉ ngón tay run run vào mặt Ryuuhou rồi lại nhìn sang bé Naomi đang nép bên chân chị gái: ``Ủa!? Khoan đã\dots\ Hai đứa nhóc này\dots\ chẳng phải là hai đứa hôm bữa đi câu cá ở cái công viên công nghệ với ông Kuon hồi đó sao!?''

Căn phòng bỗng chốc khựng lại một nhịp.

Ryuuhou giật mình chớp chớp mắt, đưa chiếc gậy bóng chày gãi gãi mái tóc tém, cất giọng đầy ngờ vực: ``Ủa? Sao chị lại biết hôm bữa bọn em đi câu cá ở công viên điện tử? Bữa đó là chủ nhật, bọn em đi chơi riêng với Onii mà, làm gì có gặp ai tên Sakamata đâu?''

Nghe đến đó, Chloe chết sững toàn tập.

Một giọt mồ hôi lạnh toát bất giác lăn dài từ thái dương cô nàng xuống gò má đỏ bừng. Trí não của cô nàng cá voi lập tức hú còi báo động đỏ: ``\textit{Chết dở rồi! Lỡ mồm để lộ bí mật nhiệm vụ trinh sát của tổ chức holoX rồi!! Làm sao mà giải thích chuyện mình lén lút trốn sau màn hình trung tâm giám sát rồi bám đuôi ông quản lý được đây!?}''

Trong cơn hoảng loạn tột độ, Chloe cuống cuồng vung cả hai tay lên trời xua lia lịa: ``À-À! Kh-Không có gì đâu! Ý chị là\dots\ là hôm đó chị tình cờ lướt thấy video clip dạo trên mạng thôi mà! Đúng rồi, mạng xã hội ấy! Hôm đó có người đi đường quay lại cảnh một cô bé tomboy giật được con cá chép to đùng ở hồ điện tử rồi đăng lên mạng viral, chị vô tình xem trúng thôi chứ bộ! Làm gì có chuyện một sát thủ dọn dẹp hàng đầu như chị lại rảnh rỗi đi bám đuôi ông anh trai lười biếng khó ưa của mấy đứa làm gì chứ, poe!!''

Thế nhưng, vì mải vung tay chối đây đẩy, vạt khăn tắm quanh ngực cô nàng bỗng trượt nhẹ xuống một phân, làm lộ ra một khoảng xương quai xanh trắng nõn nà.

``\textbf{Hyaaaa!!}''

Chloe hét toáng lên một tiếng the thé, vội vàng chụp hai tay túm chặt vạt khăn kéo giật lên tận cằm, cả người cuộn tròn lại như một con nhím xù lông, mặt mũi đỏ lựng đến mức tưởng chừng có thể bốc khói nghi ngút.

Ryuuhou đứng nhìn màn kịch giấu đầu hở đuôi vụng về đến mức trẻ con cũng nhìn thấu ấy, không kìm được nữa mà ôm bụng bật cười sặc sụa: ``Bwahaha! Chị lươn lẹo y chang lời Onii vừa kể xong luôn á! Chối thế này thì có ma nó tin! Hóa ra chị cá voi sát thủ lừng danh của holoX lại chính là cô bạn thanh mai trúc mã của Onii nhà em à? Em nghe danh chị lười tắm lừng lẫy trên mạng suốt mấy tháng nay rồi, nay mới được tận mắt chiêm ngưỡng `phiên bản thơm tho sạch bong' ngoài đời thực đấy! Trông chị đáng yêu hơn trên stream nhiều luôn!''

``Chị\dots\ chị không có ở bẩn mà!'' Chloe vùi nửa khuôn mặt đỏ gay vào lớp khăn bông, phụng phịu cãi bướng bằng chất giọng lí nhí: ``Đã bảo là người ta tiết kiệm nước để bảo vệ môi trường sinh thái biển thôi mà\dots\ Poe poe\dots''

Giữa lúc Chloe đang ngượng ngùng đến mức chỉ muốn tìm một cái khe sàn để chui xuống trốn, thì một bàn tay nhỏ xíu, mềm mại khẽ chạm vào vạt khăn tắm nơi đầu gối cô.

Bé Naomi ôm chú búp bê vải trong tay, chầm chậm bước đến bên mép ghế sô-pha. Cô bé út ngước đôi mắt màu xanh lục bảo trong veo như hạt sương sớm, chăm chú ngắm nhìn khuôn mặt đang đỏ ửng của Chloe. Rồi bằng một giọng nói trẻ thơ ngây thơ và ngọt ngào như rót mật, Naomi thỏ thẻ: ``Chị cá voi ơi\dots\ Chị thơm quá chừng luôn á. Giống hệt mùi hoa cúc dại mẹ trồng sau góc hè nhà mình vậy! Mắt của chị đỏ lấp lánh đẹp ơi là đẹp, giống hệt viên ngọc bích đỏ của nàng công chúa biển sâu trong cuốn truyện tranh mà Suzuno-oneechan hay đọc cho em nghe trước khi đi ngủ luôn!''

Lời khen ngây thơ, thuần khiết và chân thành đến tuyệt đối của một đứa trẻ bốn tuổi bỗng chốc như một mũi tên êm ái bắn thẳng vào tim cô nàng cá voi sát thủ.

Toàn bộ những gai góc, những bướng bỉnh và vẻ xù lông phòng thủ thường thấy của một thành viên holoX lập tức tan biến thành mây khói. Chloe sững sờ chớp chớp mắt. Đôi đồng tử đỏ rực khẽ lay động, rồi khóe môi cô nàng bất giác cong lên thành một nụ cười hiền dịu lạ thường. Cô khẽ buông một bàn tay khỏi khăn áo, cẩn thận vươn ngón tay chạm nhẹ lên má bánh bao phúng phính của Naomi, giọng mềm nhũn: ``C-Cảm ơn em nhé, bé con\dots\ Em ngoan và đáng yêu quá chừng à\dots''

Tôi ngồi trên ghế đối diện, nhìn thấy cảnh tượng ấy mà trong lòng không khỏi bật cười ấm áp. Tôi đặt chén trà xuống, vẫy tay gọi Ryuuhou lùi lại một chút rồi nhẹ nhàng cất lời: ``Để tôi giới thiệu đàng hoàng lại cho bà nhé, Chloe. Con bé tomboy nghịch ngợm cầm gậy bóng chày này là Hikawa Ryuuhou, em gái ruột của tôi. Hồi gia đình tôi chuyển từ làng chài lên Kawachi được đúng một tháng thì con bé mới chào đời, thế nên ngày xưa ở quê bà mới chưa từng nghe nhắc đến nó. Tính tình thì bộc trực, ham mê thể thao y chang con trai, nhưng được cái rất thẳng thắn và thương anh chị.''

Ryuuhou giơ hai ngón tay chữ V làm điệu, toe toét cười chào đàn chị. ``Ossu\~{} Hikawa Ryuuhou đây\~{}''

Tôi lại chỉ sang Naomi: ``Còn đây là bé út Naomi, công chúa nhỏ đáng yêu nhất nhà, lúc nào cũng ngoan ngoãn và nghe lời.''

Chloe gật gật đầu, ánh mắt vẫn còn dán vào gương mặt bầu bĩnh của Naomi. Rồi cô nàng ngước mắt nhìn sang Suzuno và Kaiho - hai gương mặt giống hệt hai anh em tôi, nhưng toát ra hai thần thái hoàn toàn đối lập - trong ánh mắt không giấu nổi sự thắc mắc: ``Còn hai người này\dots? Hồi nãy trên xe tôi nghe ông gọi là Kaiho với Suzuno đúng không? Hai cô em gái này trông chẳng giống nét mặt của ông hay bác Kaien một chút nào cả\dots''

Căn phòng thoáng chốc chùng xuống một nhịp tĩnh lặng nhẹ nhàng.

Kaiho vẫn đứng dựa lưng vào cạnh tủ, khoanh tay trước ngực với vẻ thản nhiên quen thuộc. Suzuno thì khẽ nghiêng đầu mỉm cười dịu dàng, ánh mắt trong trẻo không hề có lấy một gợn sóng bối rối.

Tôi nhìn Chloe, bằng một giọng nói trầm ấm và chân thành nhất, tôi chậm rãi giải thích:

``Đúng vậy. Về Kaiho, Suzuno, và cả bé Naomi nữa\dots\ Thực ra chúng tôi không cùng chung huyết thống. Nhưng giữa cuộc đời rộng lớn này, có những cuộc gặp gỡ do cơ duyên trời định sắp đặt, vượt qua mọi sự tình cờ thông thường. Mấy đứa nhỏ từng trải qua những năm tháng bão giông, gian khó và chịu nhiều mất mát, trôi dạt giữa muôn ngả cuộc đời. Rồi nhờ một nhân duyên kỳ diệu đưa đẩy, chúng tôi đã tìm thấy nhau giữa những ngã rẽ cuộc sống. Bố mẹ tôi đã đón nhận các em, và từ lâu, cả nhà đã coi nhau như anh chị em ruột thịt chung một mái nhà, chia ngọt sẻ bùi không một chút phân biệt.''

Tôi giữ trọn vẹn lời hứa bảo vệ bí mật về The Void và những thế giới song song trong lòng. Với Chloe - một người bạn thuở nhỏ vừa mới tái ngộ, một cô gái bình thường đang nỗ lực hết mình với ước mơ streamer - sự thật về những vết nứt không gian hay đa vũ trụ chỉ mang lại sự hoang mang không cần thiết. Điều quan trọng nhất, chân thực nhất mà cô nàng cần biết, chính là tình thân máu thịt gắn kết dưới mái nhà này.

Chloe lắng nghe từng lời của tôi, đôi mắt đỏ thoáng chùng xuống, ánh lên một sự rung động sâu sắc.

Là một đứa trẻ từ nhỏ đã phải chịu cảnh lẻ loi, một mình mang chiếc mặt nạ sát thủ gai góc bước lên chốn đô thành Tokyo đầy bon chen và cạm bẫy, Chloe hiểu hơn ai hết giá trị vô giá của hai chữ ``mái ấm''. Không có sự phán xét, không có sự dò xét tò mò, chỉ có một sự đồng cảm ấm áp tự nhiên len lỏi vào trái tim cô gái trẻ.

``Thì ra là vậy\dots'' Chloe khẽ gật đầu, cất giọng thì thầm: ``Một gia đình có nhiều nhân duyên kỳ diệu thật đấy\dots''

Đúng lúc ấy, Suzuno bước tới bên cạnh chiếc ghế sô-pha. Trên tay cô bé bưng chiếc khay gỗ nhỏ, bên trên đặt một chiếc máy sấy tóc loại êm, một chiếc khăn bông khô ráo mềm mại và một tách trà gừng mật ong nóng hổi đang nghi ngút khói.

Suzuno quỳ một gối xuống bên cạnh Chloe, nở một nụ cười hiền hậu tựa như thiên thần: ``Sakamata-san, tóc của chị vẫn còn ẩm ướt nhiều lắm đấy ạ. Nếu để lâu trong phòng điều hòa thế này chị sẽ rất dễ bị trúng gió và cảm lạnh. Để em sấy khô tóc giúp chị nhé? Chị uống một chút trà gừng mật ong này đi ạ, trà nóng sẽ giúp chị ấm bụng và lưu thông khí huyết sau khi tắm nước nóng lâu.''

Chloe ngơ ngác nhìn tách trà bốc khói thơm lừng mùi gừng và mật ong rừng, rồi lại nhìn nụ cười ngọt ngào, dịu dàng không tì vết của cô em gái tóc nâu. Lớp vỏ bọc phòng thủ trong lòng cô nàng bỗng chốc vỡ tan từng mảnh. Chloe đón lấy tách trà, hai mắt rơm rớm ngấn lệ, cảm động đến mức cất giọng mếu máo:

``Oaaa\dots\ Cảm ơn em nhiều lắm, Suzuno-chan! Trời đất ơi, sao em gái của tên Kuon này lại dịu dàng, nết na và chu đáo như một cô tiên thế này chứ!? Đúng là không bù cho cái tên anh trai khô khan, độc miệng, lúc nào cũng chỉ biết châm chọc người khác tí nào hết á!''

Tôi đứng bên cạnh chỉ biết nhún vai cười trừ.

Trong khi đó, Kaiho khẽ bước tới, tựa hông vào lưng ghế sô-pha ngay phía sau Chloe. Cô nàng đưa đôi mắt sắc sảo quét nhìn Chloe một lượt từ đầu đến chân theo đúng tác phong kiểm định viên chuyên nghiệp, rồi buông một câu nhận xét lạnh băng chuẩn thương hiệu Kaiho: ``Đỡ mùi hơn hẳn lúc ở văn phòng rồi đấy. Mẹ mà ra tay đúng là chuẩn không cần chỉnh, không bao giờ có sản phẩm lỗi. Nhưng mà lần sau đừng có giãy giụa như cá chép mắc cạn thế nữa nhé. Ghì bà chị mệt đứt cả hơi, suýt nữa thì làm đổ cả lọ dầu gội đầu đắt tiền của tôi.''

Chloe quay phắt lại, bĩu môi phồng má với cô nàng tóc đen: ``Biết rồi mà! Ai mượn mấy người mạnh tay thế làm chi, người ta có phải bao cát đâu chứ, poe!''

Cả căn phòng bỗng chốc bật cười rôm rả. Bầu không khí ngột ngạt, ngượng ngùng lúc ban đầu đã hoàn toàn tan biến, nhường chỗ cho sự ấm cúng và thân mật tự nhiên của một mái nhà.

\ct

\stpt{-Histoire 5-}
\stcap{Sự thấu hiểu của những người đi trước}

Giữa những tiếng cười đùa rộn rã của mấy chị em, tiếng bước chân thong dong nhưng đầy đanh thép từ hướng cửa bước vào bỗng thu hút sự chú ý của mọi người.

Kson bước tới gần chiếc bàn trà. Trên tay cô là lon nước ngọt có ga vừa uống dở, cơ bắp cánh tay săn chắc toát ra một thứ ``bá khí giang hồ'' ngùn ngụt của một cựu đại tỷ từng tung hoành khắp những nẻo đường nước Mỹ lẫn khu phố đèn đỏ Kansai.

Cô nàng dừng bước trước mặt Chloe, khoanh một tay trước ngực, nghiêng đầu nhìn cô nàng cá voi sát thủ đang quấn khăn tắm ngồi thu lu trên ghế. Đôi mắt sắc lẹm của nàng Rồng nheo lại, rồi một tràng cười sảng khoái, vang dội như chuông đồng vang lên khắp tầng ba:

``Gwahahaha! Gì đây, gì đây!? Đây là sát thủ dọn dẹp hiện trường khét tiếng của Secret Society holoX đấy à!? Kẻ mang mộng tưởng to tát muốn xâm chiếm thế giới mà rốt cuộc lại bị hai mẹ con nhà Hikawa tóm cổ xách về đây kỳ cọ tẩy uế như một đứa nhóc học sinh tiểu học thế này sao!?''

Chloe giật thót mình. Phong thái giang hồ bặm trợn, chất giọng oang oang cùng khí thế ngút trời của Kson ngay lập tức kích hoạt bản năng sợ hãi của cô nàng. Chloe hoảng hồn co rúm người lại, hai tay kéo khăn tắm trùm kín tận cổ, mắt trợn tròn nhìn Kson như thể đang đối diện với một trùm mafia thứ thiệt vừa bước ra từ phim hành động Hollywood: ``H-Hyaaa\dots! Ch-Chị là ai thế ạ\dots!? Kuon ơi, nhà ông chứa chấp cả đại ca xã hội đen ở tầng trên thật đấy à\dots!?''

Thế nhưng, ngay khi câu nói hoảng hốt của Chloe vừa dứt, Kson đã cười lớn một tiếng. Cô nàng thò tay vào chiếc túi nilon vừa mang từ tầng bảy xuống, lấy ra một gói snack khoai tây vị mù tạt wasabi cay nồng cỡ đại rồi ném thẳng vào lòng Chloe đánh *bẹp* một cái.

Chưa dừng lại ở đó, bàn tay to khỏe, ấm áp của nàng Rồng vươn tới, vò mạnh lên mái tóc vừa sấy của Chloe một cách cực kỳ bỗ bã nhưng chan chứa tình cảm: ``Nói bậy bạ gì đấy con nhóc này! Chị là Kson, cư dân tầng bảy của cái ổ chuột ấm cúng này đây! Cơ mà công nhận em thơm tho thế này trông đáng yêu, thuận mắt hơn hẳn cái đống giẻ lau bốc mùi lúc nãy đấy, con cá voi ngốc! Ở cái nhà này thì dẹp hết mấy cái chức danh sát thủ với tổ chức bí mật đi. Bị hai mẹ con nhà Hikawa bắt nạt thì cứ việc leo lên tầng bảy gõ cửa phòng chị, đại tỷ đây bảo kê cho hết!''

Chất giọng nửa Nhật nửa Anh bỗ bã, phong trần mà tràn đầy sự hào hiệp, chở che của Kson khiến Chloe hoàn toàn sững sờ. Đôi mắt đỏ của cô nàng chớp chớp nhìn gói snack cay trong lòng, rồi bất giác nở một nụ cười ngây ngô, mọi nỗi sợ hãi ban đầu bỗng chốc tan biến sạch trơn.

``V-Vâng\dots\ Em cảm ơn Kson-senpai ạ\dots''

``Thế có phải ngoan không!'' Kson cười khà khà, kéo chiếc ghế đẩu ngồi phịch xuống bên cạnh bàn.

Lúc này, một bóng dáng thanh thoát, nhẹ nhàng khác cũng khẽ bước tới.

Delutaya bưng chiếc đĩa sứ đựng những chiếc bánh quy bơ vừa nướng thơm phức cùng một ly cacao nóng hổi đặt nhẹ nhàng lên bàn trà trước mặt Chloe. Bộ váy mặc nhà màu lam nhạt tôn lên vẻ dịu dàng, trầm tĩnh của cô gái từng đi qua bão giông.

Nàng Tam giác khẽ mỉm cười, ngồi xuống chiếc ghế đơn ngay sát cạnh chiếc sô-pha của Chloe. Đôi mắt màu thạch anh tím của cô ánh lên một sự thấu hiểu sâu sắc và bao dung đến lạ kỳ. ``Em ăn thử một chiếc bánh quy nhé, Chloe-chan. Bánh này chị vừa nướng xong cách đây ít phút, vẫn còn giòn và ấm lắm đấy,'' Delutaya cất giọng trong trẻo, êm dịu như một khúc hát ru giữa buổi chiều tà.

Chloe ngước nhìn người con gái trước mặt mình. Dù chưa từng chính thức chạm mặt nhau trong giới giải trí, nhưng Chloe thừa biết người đang ngồi trước mặt mình là ai - Mano Aloe, tiền bối lừng danh thế hệ thứ năm, người từng là tâm điểm của một cơn giông bão dư luận kinh hoàng trước khi biến mất khỏi ánh đèn sân khấu để tái sinh thành Delutaya kiên cường của ngày hôm nay.

Chloe khẽ cầm lấy một chiếc bánh quy, cắn một miếng nhỏ. Vị bơ béo ngậy hòa quyện cùng vị ngọt thanh tan chảy trên đầu lưỡi, xua đi toàn bộ sự mệt mỏi sau cuộc giằng co trong phòng tắm.

``Bánh ngon lắm ạ\dots\ Cảm ơn chị Delu-senpai\dots'' Chloe lí nhí, hai vai khẽ chùng xuống.

Delutaya nhìn đôi bàn tay nhỏ nhắn của Chloe vẫn còn hơi run run cầm chiếc bánh, ánh mắt cô đượm vẻ thương cảm của một người đi trước đã từng nếm trải mọi đắng cay. Delutaya khẽ vươn tay, nhẹ nhàng nắm lấy mu bàn tay của cô nàng cá voi sát thủ.

Sự ấm áp từ bàn tay Delutaya khiến Chloe khẽ giật mình ngước lên.

``Thời gian qua\dots\ em đã phải chịu nhiều áp lực và vất vả lắm rồi đúng không, Chloe-chan?'' Delutaya cất lời, giọng nói dịu dàng đến mức khiến khóe mắt Chloe bỗng cay xè.

Là người từng bước chân vào thế giới của những buổi phát sóng trực tiếp, từng bị bủa vây bởi hàng vạn ánh mắt săm soi và những kỳ vọng nặng nề, Delutaya đọc thấu hoàn toàn những gì đang ẩn giấu sau vẻ ngoài bướng bỉnh của Chloe.

Cái vỏ bọc một sát thủ dọn dẹp ngổ ngáo, cái tật lười tắm bị người ta đem ra trêu chọc, cái điệu cười ``poe poe'' tưng tửng cố tạo vẻ vô tư ấy\dots\ thực chất chỉ là một lớp áo giáp vụng về mà một cô gái hướng nội dựng lên để bảo vệ bản thân trước sự khắc nghiệt của thế giới ảo.

Delutaya khẽ mỉm cười, ánh mắt nhìn về phía xa xăm như nhớ lại những ngày tháng tăm tối nhất của chính mình: ``Hồi mới debut, chị cũng từng sợ hãi đến mức nghẹt thở. Cảm giác mỗi khi bật chiếc camera lên, thấy hàng chục ngàn con người đang dõi theo từng cử chỉ, từng lời nói của mình\dots\ nó nặng nề như một tảng đá đè lên lồng ngực vậy. Có những đêm tắt máy tính đi rồi, ngồi một mình trong căn phòng bốn bức tường tối om, nhìn những con số subscribe tăng lên mà lòng chị lại trống rỗng và cô đơn đến phát khóc. Chị từng sợ mình làm không đủ tốt, sợ người ta ghét bỏ mình, và sợ đến mức chỉ muốn trốn tiệt vào một góc tối không bao giờ bước ra ánh sáng nữa.''

Chloe run rẩy lắng nghe, từng câu từng chữ của Delutaya như đánh trúng vào nỗi niềm sâu kín nhất mà cô chôn chặt trong lòng suốt bao tháng ngày qua kể từ khi debut cùng holoX. Đôi đồng tử đỏ rực của cô nàng ngân ngấn nước mắt.

Delutaya siết nhẹ bàn tay Chloe, ánh mắt tràn đầy sự khích lệ và tin tưởng: ``Thế nhưng em biết không, Chloe-chan? Em không hề đơn độc như em nghĩ đâu. Bước chân vào con đường này, em có những người đồng đội tuyệt vời bên cạnh, có những khán giả thực lòng yêu quý con người thật của em. Và đặc biệt nhất\dots\ em có Kuon-kun.''

Cô nàng liếc mắt nhìn sang tôi, khẽ mỉm cười: ``Cái cậu Tổng quản lý này trông thì lúc nào cũng làm mặt nghiêm, thích cằn nhằn và khó tính như một ông cụ non vậy thôi. Nhưng từ tận đáy lòng, cậu ấy là người chu đáo và tận tụy nhất mà chị từng biết. Cậu ấy chưa từng bỏ rơi bất kỳ ai trong chúng ta cả. Và cả mái nhà Hikawa này nữa, nơi đây sẽ luôn rộng mở để đón nhận em, bảo bọc em mỗi khi em cảm thấy mỏi mệt ngoài kia.''

Từng lời nói chân tình, ấm áp xuất phát từ tận đáy lòng của người đàn chị từng trải như một dòng suối ấm tưới mát tâm hồn khô cằn và cô đơn của Chloe. Tảng băng phòng thủ, sự ngượng ngùng và mặc cảm của một kẻ hướng nội trong lòng cô bé cá voi sát thủ hoàn toàn tan rã.

Một giọt nước mắt lăn dài qua gò má đỏ bừng của Chloe, rơi xuống mu bàn tay. Cô vội vã đưa vạt khăn lau vội, rồi gật đầu thật mạnh, nở một nụ cười rạng rỡ và nhẹ nhõm nhất kể từ khi bước chân vào ngôi nhà này:

``Dạ\dots! Em hiểu rồi ạ\dots\ Em cảm ơn chị Delu nhiều lắm\dots!''

Kson bên cạnh vỗ đùi đánh *ĐÉT* một cái, hào sảng giơ lon nước ngọt lên: ``Đấy! Phải thế chứ! Cười lên thế có phải xinh xắn hơn bao nhiêu không!''

Ryuuhou và bé Naomi cũng vỗ tay reo hò, tiếng cười nói rộn rã làm rực sáng cả gian phòng khách tầng ba. Hương thơm của hoa cúc, của trà gừng mật ong, của bánh quy bơ mới nướng và những tiếng cười ấm áp quấn quýt lấy nhau, dệt nên một bức tranh gia đình trọn vẹn và êm đềm đến nao lòng.

Từ phía cầu thang dẫn lên tầng sáu, tiếng xèo xèo của chảo gang cùng mùi sườn nướng mật ong ngạt ngào bắt đầu bay thoang thoảng xuống, báo hiệu cho bữa tiệc tối thịnh soạn đang sắp sửa hoàn tất.

Thế nhưng trước khi bữa tiệc bắt đầu, Chloe khẽ ngước mắt lên, len lén đưa ánh nhìn về phía tôi. Một cái nhìn đầy ngụ ý, như một lời nhắc nhở thầm lặng về một cuộc hẹn chưa trọn vẹn trên khoảng sân lộng gió của tầng thượng\dots

\ct

\stpt{-Histoire 6-}
\stcap{Tâm sự ban công lộng gió và lời ước hẹn Kawachi}

Khoảng hai mươi phút sau cuộc trò chuyện rôm rả ở phòng khách tầng ba, Chloe đã mượn tạm phòng nghỉ của Suzuno để thay sang bộ trang phục thường ngày. Đó là bộ quần áo do người quản lý phụ trách của cô nàng đã chu đáo gói ghém sẵn trong chiếc túi xách cứu trợ gửi theo từ văn phòng COVER.

Khi cánh cửa tầng thượng mở ra, gió đêm mùa hè của vùng thủ đô lập tức ùa tới, lồng lộng thổi tung tà áo và mang theo hơi thở mát rượi sau một ngày dài oi ả.

Tôi quay đầu lại và khẽ mỉm cười.

Không còn là tấm khăn tắm to sụ quấn vội hay chiếc áo khoác bốc mùi nhếch nhác thường ngày, Chloe lúc này trông gọn gàng và nữ tính đến lạ lùng. Cô nàng diện một chiếc áo váy ôm sát cơ thể với họa tiết so le hai màu trắng và xám tro. Một sự phối màu tối giản nhưng tinh tế, vừa khéo léo gợi nhắc đến tông màu đặc trưng của loài cá voi sát thủ, vừa tôn lên vóc dáng thanh mảnh, nhỏ nhắn. Mái tóc màu xám tro mềm mại đã được sấy khô hoàn toàn, buông xõa bồng bềnh sau lưng, bay nhè nhẹ theo từng đợt gió đêm. Trên mái tóc ấy, cô cài hờ một chiếc kẹp tóc nhỏ hình vây cá màu bạc.

Hương xà phòng hoa cúc dại thoang thoảng quyện cùng mùi tinh dầu tràm dịu nhẹ phảng phất trong không gian, thanh sạch và dễ chịu vô cùng.

``Cầm lấy này,'' tôi bước tới bên lan can sắt, chìa ra một lon nước ép đào ướp lạnh.

Chloe chớp chớp mắt, đưa hai tay đón lấy lon nước. Tiếng kim loại bật nắp vang lên một tiếng *TÁCH* giòn tan. Cô nàng áp mặt ngoài lành lạnh của lon nước lên gò má vẫn còn hơi ửng hồng của mình, khẽ thở hắt ra một hơi khoan khoái: ``Oa\dots\ Mát lạnh đã quá! Đúng là chỉ có nước trái cây ướp lạnh sau khi tắm xong mới là chân ái thôi!''

``Mẹ tôi mà nghe thấy câu này là bà lại xách tai lôi vào bồn tắm lần nữa đấy,'' tôi bật cười, cũng mở lon nước ép của mình rồi nhấp một ngụm ngọt lịm. ``Bà vừa khen tắm xong sảng khoái đấy nhé, đừng có mà tuần sau lại chứng nào tật nấy, trốn tắm cả tháng trời nữa.''

``Xí! Biết rồi, không cần ông phải cằn nhằn đâu, đồ khó tính!''

Chloe bĩu môi, phụng phịu lườm tôi một cái nhưng trong ánh mắt chẳng hề có lấy một tia giận dỗi nào. Cô nàng bước tới sát mép ban công, tựa hai cùi chỏ lên thanh lan can sắt, hai bàn tay đan vào nhau nâng cằm, chăm chú phóng tầm mắt ra khoảng không gian bao la phía trước.

Tầng thượng của khu ký túc xá Hikawa có tầm nhìn tuyệt đẹp. Từ độ cao này, toàn cảnh thành phố Kawachi về đêm hiện lên như một tấm thảm nhung rực rỡ, lung linh ánh đèn vàng cam ấm áp trải dài từ chân đồi ra tận phía chân trời. Xa xa, đoàn tàu điện ngầm tuyến Kintetsu lướt êm ả như một dải lụa phát sáng cắt ngang màn đêm, để lại những âm thanh xình xịch trầm đục hòa cùng tiếng gió rít qua rặng phi lao.

Lần đầu tiên sau mười sáu năm đằng đẵng, giữa chúng tôi mới có một khoảng không gian thực sự riêng tư và tĩnh lặng đến thế này. Không còn là mối quan hệ giữa một vị Tổng quản lý quyền lực, nghiêm khắc với một cô nàng thực tập sinh bướng bỉnh trong văn phòng công ty. Ở đây, dưới bầu trời đêm lộng gió của Kawachi, chúng tôi chỉ đơn thuần là Kuon và Chloe - hai đứa trẻ của làng chài ven biển năm nào.

``Này, Kuon\dots'' nàng Cá voi khẽ cất tiếng, giọng nói cô nàng nhỏ nhẹ, êm ả lạ thường, gần như chìm vào tiếng gió đêm.

``Ừ?''

``Cái chức danh `Sát thủ dọn dẹp hiện trường' của Secret Society holoX ấy\dots\ Nghe ngầu lắm đúng không?''

Tôi nghiêng đầu nhìn cô nàng, khóe môi khẽ cong lên một nụ cười trêu chọc: ``Nghe thì sát khí đằng đằng đấy chứ. Hồ sơ nhân vật bảo là chuyên gia trừ khử nhân chứng, dọn dẹp mọi dấu vết tội phạm trong bóng tối, nghe cứ như một nhân vật phản diện nguy hiểm bậc nhất trong phim hành động vậy.''

``Xạo ke hết đấy!'' Chloe thở dài thượt một hơi, hai vai xụi lơ tựa hẳn người vào lan can. ``Sát thủ cái nỗi gì chứ! Thực chất trong cái hội holoX dở hơi đó, tôi chỉ là một đứa thực tập sinh kiêm lao công quét rác đúng nghĩa đen thì có, poe!''

Tôi bật cười thành tiếng. Quả đúng là như vậy. Trong cái tổ chức bí mật tự xưng là muốn ``xâm chiếm thế giới'' ấy, cô nàng cá voi sát thủ này luôn là người phải chịu nhiều ấm ức nhất.

``Ông không biết đâu, ngày nào ở văn phòng tôi cũng muốn phát điên lên được!'' Chloe bắt đầu giãi bày nỗi lòng với gương mặt méo mó dở khóc dở cười. ``Nào là con bé tổng soái La+ hở một tí là ăn vạ, đập phá đồ đạc đòi ăn bánh ngọt; rồi cô nàng bác sĩ điên Koyori tối ngày pha chế mấy cái ống nghiệm phát nổ làm khói đen bốc nghi ngút khắp phòng thí nghiệm; Iroha thì vung kiếm chém bay cả rèm cửa; còn chị Lui thì suốt ngày dí deadline dí giấy tờ! Thế là ai phải đi hốt rác, lau sàn, thu dọn tàn cuộc sau mấy trò quậy phá long trời lở đất đó hả!? Là tôi chứ ai!! Một đứa đến cả chuyện tắm rửa cho bản thân còn lười như tôi mà ngày nào cũng phải cầm cây lau nhà với xô nước đi dọn dẹp hiện trường cho lũ dở hơi ấy, nghĩ có ức chế không cơ chứ!?''

Nhìn bộ dạng vung tay múa chân kể khổ đầy sống động của Chloe, tôi không nhịn được mà bật cười khúc khích. Nhưng khi tràng cười dứt điểm, tôi lại cảm nhận được một nỗi niềm sâu lắng hơn đang ẩn giấu sau những lời than thở nửa đùa nửa thật ấy.

Tôi hiểu. Cái mác sát thủ dọn dẹp gai góc, lạnh lùng, cái tính lười nhác và điệu cười tưng tửng trên sóng stream\dots\ thực chất chỉ là một lớp vỏ bọc mà cô nàng vụng về tạo ra.

Chloe khẽ nhấp một ngụm nước ép đào, ánh mắt lại hướng về phía những dải đèn đêm lấp lánh dưới phố, giọng nói bỗng chùng xuống, đượm một nỗi buồn man mác: ``Thực ra\dots\ tôi đâu có muốn làm sát thủ hung dữ gì đâu. Từ nhỏ, tôi chỉ thích hát thôi. Tôi chỉ muốn được đứng trên sân khấu, hát những bài hát mình yêu thích, mang lại niềm vui cho mọi người và được mọi người yêu quý, công nhận\dots\ Nhưng mà, bước chân vào giới này rồi mới thấy nó đáng sợ đến nhường nào.''

Cô nàng siết chặt lon nước trong tay, các đốt ngón tay hơi trắng bệch.

``Người ta soi mói từng lời nói, từng cử chỉ nhỏ nhặt nhất. Một cô gái vụng về, ít nói, hướng nội sâu sắc như tôi\dots\ sợ bị người ta phán xét, sợ bị ghét bỏ đến phát khóc. Thế nên tôi mới nghĩ, cứ khoác lên mình một chiếc mặt nạ gai góc, bất cần, giả vờ như mình là kẻ ngổ ngáo chẳng sợ trời chẳng sợ đất, thì sẽ không ai có thể làm tổn thương mình được nữa\dots\ Nhưng mà mệt mỏi lắm, Kuon à. Nhiều đêm về đến căn phòng trọ tối om ở Tokyo, tôi thấy ngột ngạt và cô đơn đến mức chẳng buồn bước chân vào nhà tắm nữa. Cứ thế mà gục đầu xuống ngủ thiếp đi trong bóng tối\dots''

Gió đêm thổi qua làm những lọn tóc xám của Chloe bay lòa xòa trước trán. Cô nàng quay sang nhìn tôi, đôi mắt đỏ rực long lanh ngấn lệ, bờ môi khẽ run run như muốn hờn dỗi.

``Hồi mười sáu năm trước\dots\ ngày gia đình ông chuyển đi, ông có biết tôi đã khóc nhiều thế nào ở bến xe buýt không? Bãi cào ngao ven biển, bờ kè đá nhảy, hàng rào gỗ thấp trước sân nhà\dots\ chỗ nào cũng chỉ còn lại một mình tôi. Thế mà ông đi biền biệt, không thèm gửi về lấy một cánh thư, cũng chẳng có một cuộc điện thoại nào cả. Suốt bao nhiêu năm trời, tôi cứ ngỡ\dots\ tôi cứ ngỡ ông đã quên sạch con bé cá voi ngốc nghếch này rồi chứ\dots''

Một giọt nước mắt trong veo khẽ lăn dài trên gò má trắng ngần của Chloe, lấp lánh dưới ánh trăng mờ.

Trái tim tôi bất giác nhói lên một nhịp. Tôi đặt lon nước ép xuống bờ lan can, ngước nhìn lên vòm trời đêm đầy sao của Kawachi, cất giọng trầm ấm, chân thành từ tận đáy lòng: ``Làm sao tôi có thể quên được chứ, Sakamata.''

Chloe hơi ngẩng đầu lên, ngơ ngác nhìn tôi.

``Rời khỏi làng chài năm mười tuổi, cuộc sống của tôi cũng chẳng hề dễ dàng như bà nghĩ đâu,'' tôi chậm rãi tâm sự, ánh mắt xa xăm. ``Chuyển đến một vùng đất mới, bố mẹ phải vất vả gầy dựng lại xưởng cơ khí và kho hàng từ hai bàn tay trắng. Rồi biến cố ập đến, việc chăm lo cho Ryuuhou, việc tìm kiếm và đón nhận những đứa em lưu lạc về dưới một mái nhà\dots\ Cuộc sống mưu sinh và trách nhiệm cứ thế cuốn tôi đi như một cơn lốc. Rồi khi bước chân vào Cover, gánh trên vai chiếc ghế Tổng quản lý hoàn toàn đột xuất, tôi phải đối mặt với muôn vàn sóng gió thương trường để bảo vệ công ty, bảo vệ các thế hệ talent khỏi những cạm bẫy của thế giới giải trí khắc nghiệt này. Có những lúc tôi tưởng chừng như mình không còn một giây phút nào cho bản thân nữa.''

Tôi quay sang nhìn thẳng vào đôi mắt đỏ của Chloe, khóe môi khẽ nở một nụ cười hiền hậu: ``Nhưng trong ký ức của tôi, hình ảnh một cô bé cá voi nhỏ thó ngậm kẹo mút dâu tây, nắm chặt tay tôi kéo ra bãi bồi cào ngao, bắt còng gió và nhảy sóng dưới ánh nắng hè rực rỡ năm nào\dots\ chưa bao giờ phai nhạt cả. Chiếc móc khóa đuôi cá voi khắc chữ `C \& K' mà bà tặng tôi ngày chia tay ở bến xe buýt, tôi vẫn luôn cất giữ cẩn thận trong ngăn kéo bàn làm việc suốt mười sáu năm qua đấy thôi.''

Chloe khẽ giật mình, đôi mắt đỏ mở to đầy ngỡ ngàng: ``Ông\dots\ ông vẫn còn giữ nó thật sao\dots?''

``Tất nhiên rồi,'' tôi mỉm cười, đưa tay cốc nhẹ một cái vào trán cô nàng như thói quen thuở nhỏ. ``Thế nên, đừng bao giờ nói những lời ngốc nghếch như vậy nữa. Cho dù bây giờ bà có khoác lên mình chiếc mặt nạ sát thủ gai góc của holoX, hay có trở thành một nữ VTuber triệu sub nổi tiếng lẫy lừng khắp thế giới đi chăng nữa\dots\ thì đối với tôi, bà vẫn mãi là con bé Sakamata Chloe ngốc nghếch, vừa lười tắm vừa sợ nước nóng của làng chài năm nào thôi.''

Câu nói nửa trêu chọc nửa chân tình của tôi khiến Chloe không kìm được nữa. Cô nàng vừa bật cười vừa thút thít đưa tay áo lau vội khóe mắt, hai má đỏ bừng lên vì xúc động và thẹn thùng: ``Cái tên đáng ghét này\dots\ Lúc nào cũng biết cách làm người ta muốn khóc hết á\dots!''

Tôi mỉm cười cầm lại lon nước ép đào, giơ lên trước mặt cô nàng: ``Nào, làm một ly chứ?''

Chloe chớp chớp mắt, rồi khóe môi cô nàng nở một nụ cười rạng rỡ, trong trẻo như ánh mặt trời rọi xuống mặt biển mùa hè. Cô nàng giơ lon nước của mình lên, chạm nhẹ vào lon nước của tôi.

*KENG!*

Âm thanh kim loại chạm nhau vang lên một tiếng giòn tan, thanh thoát giữa bầu trời đêm lộng gió của Kawachi.

Đó không chỉ là một cái chạm ly đơn thuần. Đó là một lời giao ước ngầm thiêng liêng giữa hai tâm hồn:

Ở công ty, trước mặt mọi người và khán giả, cô vẫn sẽ là cô nàng thực tập sinh holoX tinh quái, ngổ ngáo, còn tôi vẫn sẽ là ngài Tổng quản lý nghiêm khắc, kỷ luật giữ vững trật tự. Thế nhưng sau cánh gà, ngoài đời thực, tôi sẽ luôn là điểm tựa vững chãi nhất, là người bạn thuở thiếu thời sẵn sàng dang rộng đôi tay che chở cho cô bạn nhỏ, để cô có thể yên tâm sải cánh và tỏa sáng rực rỡ nhất trên sân khấu của ước mơ.

Chloe nhấp một ngụm nước ngọt lịm, thở phào nhẹ nhõm, tựa cằm vào cánh tay tôi, khẽ thì thầm: ``Cảm ơn ông nhé, Kuon\dots\ Có ông ở đây, tôi không còn cảm thấy sợ hãi nữa rồi\dots''

Thế nhưng, đúng vào khoảnh khắc êm đềm, lãng mạn và lắng đọng ấy\dots

Từ khoảng giếng trời thông gió của tầng sáu, một giọng thét oanh vàng, đanh thép và đầy uy lực bỗng xuyên thủng màn đêm, vang dội lên tận sân thượng: ``\textbf{Kuon!! Chloe!! Hai đứa trốn lên sân thượng hóng gió đủ chưa hả!? Mau xuống ăn cơm sườn nướng mật ong nhanh lên kẻo nguội ngắt bây giờ!!}''

Tiếng thét uy quyền của mẹ tôi khiến cả hai chúng tôi giật bắn mình suýt đánh rơi cả lon nước ngọt.

Chloe ngơ ngác mất một giây rồi lập tức ôm bụng bật cười ngặt nghẽo: ``Pwahaha! Bác Kaien uy lực dũng mãnh y chang hồi xưa luôn!''

``Chạy mau thôi bà tướng!'' Tôi dở khóc dở cười, vội vàng kéo tay cô nàng hướng về phía cửa cầu thang.
``Chạy thôiii! Poe poeee!''

Chloe cười giòn tan, bàn tay nhỏ nhắn nắm chặt lấy tay tôi, hai đứa cùng nhau lao nhanh xuống những bậc thang sáng đèn. Hương thơm ngào ngạt của món sườn nướng mật ong và những tiếng cười nói rôm rả từ tầng sáu đang chờ đón chúng tôi phía dưới.

Một đêm hè thật dài, nhưng tràn ngập sự ấm áp của tình thân và những mảnh ký ức ngọt ngào đã được hàn gắn trọn vẹn dưới bầu trời Kawachi.

%==========================================TRUYỆN PHỤ=========================================
\omk

Bữa tiệc sườn nướng mật ong ở gian bếp lớn tầng sáu vừa kết thúc trong sự no nê và thỏa mãn của tất cả mọi người.

Bàn ăn dài bằng gỗ sồi đã được dọn dẹp sạch bóng, thay vào đó là một đĩa dưa hấu ướp lạnh cắt lát đỏ au ngọt lịm cùng những tách trà lúa mạch thanh mát nghi ngút hương thơm. Đại gia đình mười người quây quần đông đủ quanh chiếc bàn lớn - tôi, mẹ, bốn đứa em gái ruột thịt và kết nghĩa, Kson, Delutaya, và cô nàng cá voi sát thủ Chloe đang vui vẻ gặm miếng dưa hấu to tướng với đôi má phúng phính.

Đúng lúc ấy, chiếc điện thoại thông minh của tôi đặt trên mặt bàn bỗng đổ chuông FaceTime dồn dập. Màn hình nhấp nháy hiện lên hình ảnh đại diện của Marine.

Tôi khẽ nhướng mày, với tay bấm nút nhận cuộc gọi rồi tiện tay bật luôn loa ngoài và camera trước: ``Alo, Marine à? Em về đến nhà an toàn rồi chứ?''

Vừa kết nối, hình ảnh của nàng Thuyền trưởng lập tức choán ngợp màn hình. Nàng Thuyền trưởng lúc này đã thay sang bộ đồ ngủ sọc hải tặc thoải mái, thế nhưng trên mặt em ấy vẫn cẩn thận đeo\dots\ tới hai lớp khẩu trang y tế bịt kín mít từ sống mũi xuống tận cằm.

Marine ghé sát mặt vào màn hình, cất giọng the the đầy vẻ lo lắng và thảng thốt: ``Alo alo! Kuon-senpai đấy à!? Tình hình chiến sự bên đó thế nào rồi anh!? `Kiện hàng sinh học nguy hiểm cấp độ ba' đã bị khống chế chưa? Bác Kaien và mấy đứa nhỏ có ai bị ngộ độc khí thải hay phải gọi xe cấp cứu rửa ruột không thế!? Em thề là từ chiều đến giờ cái văn phòng bên này vẫn còn ám ảnh mùi tanh tưởi, em phải đốt tới ba nén trầm hương để tẩy uế đấy ạ!''

Tôi phì cười, không nói không rằng liền xoay ngược camera một góc một trăm tám mươi độ, chĩa thẳng ống kính vào gương mặt của Chloe đang ngồi đối diện.

Nàng cá voi sát thủ lúc này đang diện chiếc váy ôm sát so le trắng-xám tinh tươm, mái tóc xám tro mềm mại buông xõa dài bồng bềnh, da dẻ trắng hồng mịn màng sạch bong kin kít. Toàn thân cô nàng tỏa ra làn hương hoa cúc dại thanh khiết, ngào ngạt đến mức bất kỳ ai ngồi gần cũng cảm thấy khoan khoái lạ thường.

Nhìn thấy khuôn mặt quen thuộc của bà chị Thuyền trưởng trên màn hình, Chloe lập tức phồng mang trợn má, nuốt vội miếng dưa hấu rồi hét toáng vào ống kính: ``Marine-senpai!! Chị nhìn cho thật kỹ vào đây này!! Em đây này!! Em đã được kỳ cọ sạch bong kin kít, thơm tho ngào ngạt từ chân tơ kẽ tóc rồi nhé!! Chị đừng có mà bôi bác em trên stream nữa đấy, đồ Thuyền trưởng già độc mồm độc miệng, poe!!''

Marine bên kia đầu dây chết lặng toàn tập.

Đôi mắt một mí của cô nàng trợn tròn lên như hai chiếc đĩa tây, chiếc khẩu trang y tế tụt hẳn xuống cổ, miệng há hốc kinh ngạc đến mức đánh rơi cả chiếc gối ôm hình cá voi: ``Ái chà chà\dots\ Sakamata đấy ư!? Thần tích giáng trần rồi!! Từ một bãi rác sinh học di động mà biến thành một nàng tiên cá thơm phức, trắng trẻo thế kia á!? Kaien-san ơi, bác đúng là vị thánh sống cứu rỗi nhân loại, là kỳ tích vĩ đại của nền công nghiệp tẩy rửa Nhật Bản rồi ạ!!''

Mẹ từ trong bếp bưng thêm đĩa hoa quả ra, nghe thấy thế liền cười khà khà, vẫy tay đầy tự đắc: ``Có gì đâu cháu! Tay nghề của bác thì cháu cứ yên tâm tuyệt đối! Lần sau con bé cá voi này mà dám trốn tắm quá ba ngày ở văn phòng, cháu cứ gọi một cú điện thoại, bác lại lái xe bảy chỗ lên tận sào huyệt hốt nó về đây kỳ cọ tiếp!''

Chloe nghe thấy thế thì rùng mình ớn lạnh từ đỉnh đầu xuống gót chân, vội vàng chắp hai tay lạy lục: ``Dạ thôi\dots\ cháu lạy bác, cháu xin chừa đến già rồi ạ\dots''

Cả căn phòng nổ ra một tràng cười nghiêng ngả. Marine vừa cười ngặt nghẽo vừa chúc cả nhà ngon miệng rồi mới yên tâm cúp máy đi ngủ.

Sau màn chào hỏi ồn ào của nàng Thuyền trưởng, câu chuyện bên bàn trà lại quay trở về với chủ đề sinh hoạt thường nhật.

Ryuuhou vừa nhai dưa hấu nhồm nhoàm vừa huých cùi chỏ vào tay chị Hai: ``Này Kaiho-nee, nghe nói chị được nghỉ phép nguyên một tuần luôn hả?''

Kaiho khoanh hai tay trước ngực, thong thả nhấp một ngụm trà lúa mạch mát lạnh, khóe môi khẽ nhếch lên một nụ cười đắc thắng hiếm thấy: ``Ừ. Chiều hôm kia ở căn mancave của tụi \hls, vừa vật tay hạ gục một mạch cả năm anh chàng bên đó\footnote{{\humanist FOLGE 23}, quyển số Ba thuộc \textit{stars*collection}}. Chị đã xin phép bên công ty nghỉ trọn vẹn một tuần tới để dưỡng sức chuẩn bị cho chuyến đi biển mùa hè sắp tới của nhà mình.''

Suzuno khẽ mỉm cười dịu dàng, lật mở cuốn sổ tay nhỏ bìa da đã ghi chép cẩn thận từng mốc thời gian: ``Lịch trình chuyến đi biển của gia đình mình đã được chuẩn bị xong xuôi rồi ạ. Theo đúng kế hoạch, cả nhà chúng ta sẽ bắt đầu khởi hành vào đêm nay. Chuyến đi sẽ kéo dài ba ngày hai đêm tại bãi biển phía nam.''

Ryuuhou hào hứng giơ nắm đấm lên trời huýt gió: ``Tuyệt cú mèo! Đồ bơi em với Suzu-nee đã đi sắm từ mấy hôm trước rồi, em còn tậu cả quả kính lặn với ống thở xịn đét nữa cơ, phen này tha hồ mà săn bạch tuộc!''

Bé Naomi ôm chặt chú búp bê vải trong lòng, tròn xoe đôi mắt màu lục bảo reo lên trong trẻo: ``Đi biển! Đi biển bắt ốc với Onii-chan và các chị ạ!''

Chloe ngồi bên cạnh, nghe đến hai chữ ``đi biển'' thì đôi tai cá voi vô hình dường như dỏng đứng lên. Cô nàng chớp chớp mắt tò mò: ``Ủa, chuyến đi biển của cả nhà đông đủ thế này\dots\ bác Ryujin có đi cùng không ạ?''

Tôi lắc đầu, chậm rãi giải thích: ``Bố anh không đi được. Ông ấy phải ở lại Kawachi vì ba lý do: một là để chăm sóc bà nội đã lớn tuổi ở tầng bốn, một mình bác Hijaki làm không xuể; hai là xưởng cơ khí dưới ngoại ô đang có mấy đơn hàng chế tạo máy đột xuất cần bố trực tiếp giám sát; và ba là lý do quan trọng nhất\dots''

Tôi khẽ đưa ngón tay chỉ ngược lên trần nhà: ``Phải có người ở lại ký túc xá để trông coi và `trấn áp' Mike-chan ở tầng tám. Nếu để cô ấy ở nhà một mình mà không ai để mắt, lỡ nửa đêm cô ấy nổi giận la hét hay đập vỡ bàn thì có mà sập cả ký túc xá. Bác Hijaki không còn đủ sức khỏe để chạy lên chạy xuống mãi được đâu.''

Ryuuhou lập tức gật gù tán đồng cái rụp: ``Chuẩn luôn! Bố là người duy nhất sở hữu độ bình thản tuyệt đối để đối phó với những cơn thịnh nộ long trời lở đất của Mikeneko-san!''

Tôi tiếp tục tóm tắt danh sách thành viên: ``Thế nên, chuyến đi lần này chỉ có mẹ Kaien làm tài xế kiêm trưởng đoàn, năm anh em nhà mình, và các cô nàng Thế hệ thứ Tư thôi.''

Kson đập bàn cười khà khà đầy hào sảng: ``Chuẩn bài luôn! Chiều nay em vừa gọi điện liên lạc cho cả bốn đứa bên kia xong. Kanatan, Đùi Cừu, Towa với cả công chúa Luna nghe tin đi biển là đứa nào đứa nấy nhảy cẫng lên reo hò ầm ĩ! Phía công ty và ban quản lý bên \hll\ cũng đã chính thức cấp phép cho bọn nó nghỉ trọn ba ngày rồi. Để em chuẩn bị một thùng bia lạnh ướp đá với mấy chai soda Mỹ, đại tỷ đây sẽ làm chân pha cocktail thượng hạng kiêm bảo kê cho các em gái suốt chuyến đi!''

Delutaya khẽ mỉm cười thanh nhã, nhẹ nhàng lắc đầu tiếp lời: ``Còn em thì xin phép cả nhà cho em ở lại ký túc xá đợt này ạ. Tâm trạng của Mike-chan dạo này vẫn còn khá nhạy cảm, chỉ có bố Ryujin với bác Hijaki thì em sợ hai người đàn ông sẽ khó mà tinh tế vỗ về, chăm sóc chị ấy được. Em ở nhà vừa có người bầu bạn với Mike-chan, lại tiện đỡ đần chăm sóc bà nội luôn. Nhưng mọi người yên tâm, sáng sớm em sẽ nướng thật nhiều bánh quy bơ, làm sandwich kẹp thịt và chuẩn bị hoa quả dầm đóng hộp để cả nhà mang theo ăn nhẹ dọc đường nhé!''

Bà Kaien nghe vậy liền mỉm cười gật đầu tán thưởng: ``Được thế thì tốt quá, Delu-chan lúc nào cũng chu đáo và biết nghĩ cho mọi người. Cứ giao nhà cửa cho con với bố nó nhé!''

Không khí bàn ăn càng lúc càng rôm rả, náo nhiệt. Và ngay lúc đó, bản năng của loài sinh vật biển sâu trỗi dậy mãnh liệt trong lòng cô nàng Sakamata Chloe.

Đôi mắt đỏ rực của cô nàng sáng bừng lên như hai ngọn đèn pha ô tô. Nghĩ đến cảnh được hòa mình vào làn nước biển mát rượi, được ăn hải sản nướng ngập mồm và ở bên cạnh những con người ấm áp này suốt ba ngày, nàng Cá voi không kìm lòng được nữa. Cô nàng lập tức buông miếng dưa hấu xuống, lon ton chạy vòng qua bàn ăn, xà vào lòng bà Kaien rồi níu chặt lấy cánh tay tôi, giở giọng làm nũng ngọt ngào như rót mật vào tai: ``Bác Kaien ơi! Kuon ơi! Cho cháu\dots\ cho cháu đi biển cùng với cả nhà với được không ạ!? Cháu là cá voi mà, lâu lắm rồi cháu chưa được bơi lội dưới biển xanh! Cháu hứa sẽ cực kỳ ngoan ngoãn, ngày nào cũng tắm rửa thơm tho, không quậy phá tí nào đâu ạ, poe poe!''

Bà Kaien bật cười sảng khoái, xoa mạnh đầu cô nàng: ``Được chứ sao không! Đi đông càng vui! Bác đang lo mấy đứa con gái nhà này lười vận động dưới nước, có thêm con cá voi sát thủ như cháu thì bãi biển nó mới náo nhiệt!''

Thế nhưng tôi lập tức nhướng mày, khoanh tay nhìn cô nàng bằng ánh mắt ngờ vực: ``Khoan đã nào, Chloe. Bà là talent chính thức của COVER, lại là thành viên của tổ chức Secret Society holoX. Lịch trình stream, luyện tập vũ đạo với họp hành của tổ chức bà vứt đi đâu? Bà tự ý biến mất ba ngày hai đêm thì sếp của bà không lột da bà ra mới lạ đấy.''

Chloe khẽ nhếch mép, nở một nụ cười ranh mãnh đầy tự tin: ``Hừm! Chuyện đó thì ông không phải lo! Để sát thủ dọn dẹp kiêm chuyên gia tình báo Sakamata Chloe này trổ tài cho mà xem!''

Nói là làm, cô nàng thoăn thoắt rút chiếc điện thoại thông minh ra, mười ngón tay lướt như múa trên bàn phím cảm ứng, soạn ngay một bản báo cáo khẩn cấp gửi thẳng vào nhóm chat nội bộ của ban lãnh đạo holoX:

\begin{quote}
    \textit{`[BÁO CÁO NHIỆM VỤ ĐẶC BIỆT CỦA SÁT THỦ DỌN DẸP SAKAMATA CHLOE]: Kính gửi Tổng soái La+ Darknesss vĩ đại! Tôi đã xuất sắc thâm nhập thành công vào sào huyệt của mục tiêu số một - Tổng quản lý Hikawa Kuon! Nhận thấy mục tiêu chuẩn bị tổ chức một chuyến dã ngoại bí mật ra vùng biển phía nam, để theo dõi trực tiếp đối tượng hai mươi tư trên bảy và thu thập các dữ liệu tình báo cơ mật nhất phục vụ cho kế hoạch xâm chiếm thế giới của tổ chức, tôi xin phép được trú lại nhà mục tiêu vài ngày và tháp tùng theo dõi chuyến đi! Yêu cầu Tổng soái phê duyệt khẩn cấp!'}
\end{quote}

Chỉ vỏn vẹn ba giây sau, chiếc điện thoại trên tay Chloe bỗng rung lên bần bật.

Một thông báo phản hồi với dòng chữ in hoa to đùng từ tài khoản của La+ Darknesss hiện rõ mồn một trên màn hình:

\begin{quote}
    \textit{`[PHÊ DUYỆT NGAY VÀ LUÔN! Làm tốt lắm Chloe! Cứ bám sát tên Tổng quản lý khó ưa đó, nhớ chụp trộm vài tấm hình dìm hàng lúc hắn mặc đồ bơi làm bằng chứng tống tiền cho ta! À mà tiện thể đi biển nhớ mua quà lưu niệm với đặc sản bánh mực nướng đem về nộp cho ta đấy!!]'}
\end{quote}

Chloe giơ thẳng chiếc màn hình điện thoại ra trước mặt tôi, ngẩng cao cằm làm điệu bộ đắc thắng vô cùng vênh váo: ``Thấy chưa hả tên Kuon kia! Sếp của tôi đã chính thức phê duyệt nhiệm vụ gián điệp tối mật rồi nhé! Từ giờ cho đến hết tuần sau, tôi chính thức `trú đóng' tại ký túc xá này để giám sát ông toàn thời gian, chuẩn bị tinh thần phục vụ tôi đi, đồ lười biếng!''

Từ trên chiếc đồng hồ đeo tay của tôi, trợ lý Game phát ra một tiếng thở dài bằng giọng kim loại khô khốc: ``Xác suất thành công của chiến dịch tình báo này đối với kế hoạch xâm chiếm thế giới là đúng không phần trăm. Thực chất, Sakamata-san chỉ là một đứa trẻ ham ăn ham chơi thèm đi tắm biển và không muốn về nhà một mình thôi.''

Cả phòng ăn tầng sáu lập tức nổ ra một tràng cười nghiêng ngả. Ryuuhou đập bàn cười sặc sụa, Naomi ôm búp bê vỗ tay nhảy cẫng lên, còn Kson thì nâng lon nước ngọt huýt sáo vang dội cả một góc trời đêm.

Tôi lắc đầu cười trừ, nhìn cô bạn thuở nhỏ đang háo hức cùng Suzuno và Ryuuhou chụm đầu bàn tán về các loại kem chống nắng và mẫu đồ bơi mới.

Tiếng cười rộn rã của đại gia đình vang vọng khắp gian phòng ăn ấm áp, hòa cùng làn gió đêm mùa hè lồng lộng của thành phố Kawachi. Chuyến hành trình ba ngày hai đêm hướng về vùng biển đầy nắng vàng, cát trắng và những con sóng dạt dào đã chính thức được ấn định ngày khởi hành. Những trận đấu nảy lửa, những cuộc phiêu lưu bất ngờ và vô vàn tình huống dở khóc dở cười giữa các thành viên nhà Hikawa, đại diện \textit{holoForce} cùng cô nàng cá voi sát thủ đang chờ đón họ phía trước\dots

\il

\begin{center}
    \textit{(Hành trình mùa hè sôi động của đại gia đình Hikawa sẽ được tiếp nối trọn vẹn trong \textbf{Quyển B: Bãi Sóng Cồn Sôi Động} — kính mời quý độc giả cùng đón đọc!)}
\end{center}

\end{document}